\frfilename{mti.tex}
\versiondate{31.12.10/4.9.12}
\copyrightdate{2010}










\Loadfont{\twelvebf=cmbx12}

\def\chaptername{Indeks}
\def\sectionname{Topik dan hasil utama}
\def\hc{\mathop{\text{hc}}\nolimits}
\def\Krein{Kre\v{\i}n}
\def\med{\mathop{\text{med}}\nolimits}
\def\varinnerprod#1#2{#1\dotproduct#2}
\def\Fn{\mathop{\text{Fn}}\nolimits}

\volumeno=1
\volumeno=2
\volumeno=3
\volumeno=4
\volumeno=5
\ifnum\luluvolumeno>0\volumeno=\luluvolumeno\fi

\ifnum\volumeno<5\def\vfive#1{}\else\def\vfive#1{#1}\fi
\ifnum\volumeno<4\def\vfour#1{}\else\def\vfour#1{#1}\fi
\ifnum\volumeno<3\def\vthree#1{}\else\def\vthree#1{#1}\fi
\ifnum\volumeno<2\def\vtwo#1{}\else\def\vtwo#1{#1}\fi

\def\ind{\mathop{\text{ind}}}

\def\indexiiheader#1{\ifnum\volumeno<2{}\else\indexheader{#1}\fi}
\def\indexiiiheader#1{\ifnum\volumeno<3{}\else\indexheader{#1}\fi}
\def\indexivheader#1{\ifnum\volumeno<4{}\else\indexheader{#1}\fi}
\def\indexvheader#1{\ifnum\volumeno<5{}\else\indexheader{#1}\fi}

\indexmedskip


\centerline{\twelvebf Indeks Jilid 1 dan 2}




\indexmedskip






\gdef\bottomparagraph{}
\noindent{\bf\sectionname}
   \smallskip
   \let\headlinesectionname=\sectionname

Indeks umum di bawah ini dimaksudkan untuk mencakup seluruh materi.
Tak terhindarkan, banyaknya entri kadang membuat indeks tersebut kurang membantu.
Karena itu saya menyiapkan indeks yang lebih ringkas dan beranotasi lebih baik;
saya berharap indeks ini membantu pembaca memusatkan perhatian pada bidang
tertentu. Indeks ringkas ini tidak mencantumkan definisi, sebab entri bercetak
tebal dalam indeks utama seharusnya mengarahkan pembaca ke definisi secara
efisien; jika pencarian di sini tidak membuahkan hasil, tentu carilah kembali
dalam indeks utama. Entri yang berbentuk pernyataan matematika sering
menghilangkan hipotesis penting dan harus diperiksa terhadap pernyataan formal
dalam isi buku.

\indexmedskip

\indexiiheader{kontinu mutlak} fungsi bernilai real yang kontinu mutlak \S225

----- sebagai integral tak tentu 225E

fungsional aditif kontinu mutlak \S232

----- karakterisasi 232B
































































\indexiiheader{tanpa atom} ruang ukur tanpa atom

----- mempunyai unsur dengan setiap nilai ukuran yang mungkin 215D

















\indexmedskip















































































\indexheader{Borel}











himpunan Borel dalam $\BbbR^r$ 111G

----- dan ukuran Lebesgue 114G, 115G, 134F

\indexiiheader{terbatas} fungsi bernilai real bervariasi terbatas \S224

----- sebagai selisih fungsi monoton 224D

----- integral turunannya 224I

----- dekomposisi Lebesgue 226C

























\indexiiheader{Brunn} teorema Brunn--Minkowski 266C

\indexmedskip

\indexheader{Cantor} himpunan dan fungsi Cantor 134G, 134H




\indexheader{\Caratheodory} Konstruksi ukuran dari ukuran luar oleh \Caratheodory, 113C

















\indexiiheader{Carleson} Teorema Carleson (deret Fourier fungsi terintegralkan-kuadrat konvergen hampir di mana-mana) \S286








\indexheader{pusat} Teorema Limit Pusat (jumlah variabel acak bebas berdistribusi mendekati normal) \S274


----- syarat Lindeberg 274F, 274G

\indexheader{penggantian} penggantian ukuran dalam integral 235K

penggantian variabel dalam integral \S235

----- $\int J\times g\phi\,d\mu=\int g\,d\nu$ 235A, 235E, 235J

----- menentukan $J$ 235M; 

----- ----- $J=|\det T|$ untuk operator linear $T$ 263A; $J=|\det\phi'|$ untuk operator terdiferensialkan $\phi$ 263D


----- ----- ----- bila ukurannya adalah ukuran Hausdorff 265B, 265E


----- bila $\phi$ \imp\ 235G

----- $\int g\phi\,d\mu=\int J\times g\,d\nu$ 235R

\indexheader{karakteristik} fungsi karakteristik suatu distribusi probabilitas \S285

----- barisan distribusi konvergen dalam topologi samar jika dan hanya jika fungsi karakteristiknya konvergen titik demi titik 285L















































\indexheader{lengkap} ruang ukur lengkap \S212

\indexheader{pelengkapan} pelengkapan suatu ukuran 212C {\it dst.}




\indexheader{konsentrasi} konsentrasi ukuran 264H


\indexheader{bersyarat} ekspektasi bersyarat

----- dari suatu fungsi \S233

----- sebagai operator pada $L^1(\mu)$ 242J

\indexheader{konstruksi} konstruksi ukuran

----- ukuran citra 234C

----- dari ukuran luar (metode \Caratheodory) 113C

----- ukuran subruang 131A, 214A

----- ukuran produk 251C, 251F, 251W, 254C

















----- sebagai tarikan balik 234F


















\indexheader{konvergensi} teorema konvergensi (B.Levi, Fatou, Lebesgue) \S123

konvergensi dalam ukuran (topologi ruang linear pada $L^0$)

----- pada $L^0(\mu)$ \S245



----- bila Hausdorff/lengkap/dapat dimetrikkan 245E






\indexheader{konveks} fungsi konveks 233G {\it dst.}




\indexheader{konvolusi} konvolusi fungsi

----- (pada $\BbbR^r$) \S255




----- $\int h\times(f*g)=\int h(x+y)f(x)g(y)dxdy$ 255G


----- $f*(g*h)=(f*g)*h$ 255J


konvolusi ukuran

----- (pada $\BbbR^r$) \S257



----- $\int h\,d(\nu_1*\nu_2)=\int h(x+y)\nu_1(dx)\nu_2(dy)$
257B


----- dari ukuran kontinu mutlak 257F




\indexheader{terhitung}


himpunan terhitung 111F, 1A1C {\it et seq.}

ukuran terhitung-koterhitung 211R








\indexheader{pencacahan} ukuran cacah 112Bd










\indexmedskip

\indexheader{Dedekind}






\indexiiheader{Denjoy} teorema Denjoy--Young--Saks 222L




\indexiiheader{terdiferensialkan} fungsi terdiferensialkan (dari $\BbbR^r$ ke $\BbbR^s$) \S262, \S263


\indexheader{langsung} jumlah langsung dari ruang ukur 214L










\indexheader{distribusi} distribusi suatu keluarga berhingga variabel acak \S271

----- sebagai ukuran Radon 271B

----- dari $\phi(\pmb{X})$ 271J

----- dari suatu keluarga bebas 272G

----- ditentukan oleh fungsi karakteristik 285M
















\indexheader{Doob} Teorema Konvergensi Martingal Doob 275G



















\indexmedskip































\indexheader{penghabisan} prinsip penghabisan 215A

\indexheader{diperluas} garis real diperluas \S135

\indexiiheader{perluasan} perluasan ukuran 214P













\indexmedskip

\indexheader{Fatou} Lema Fatou ($\int\liminf\le\liminf\int$ untuk barisan fungsi tak negatif) 123B




\indexheader{Fej\'er} Jumlah Fej\'er (rata-rata berjalan dari jumlah Fourier) konvergen ke rata-rata lokal $f$ 282H


----- secara seragam jika $f$ kontinu 282G






















\indexheader{Fourier} deret Fourier \S282

----- konvergen dalam norma di $L^2$ 282J

----- konvergen pada titik terdiferensialkan 282L

----- konvergen ke titik tengah lompatan jika $f$ bervariasi terbatas 282O


----- dan konvolusi 282Q

----- konvergen hampir di mana-mana\ untuk fungsi terintegralkan-kuadrat 286V

\wheader{Fourier}{0}{0}{0}{36pt} transformasi Fourier

----- pada $\Bbb R$ \S283, \S284



----- rumus untuk $\int_c^df$ dalam istilah $\varhatf$ 283F

----- dan konvolusi dari fungsi 283M




----- melalui aksinya pada fungsi uji 284H {\it dst.}

----- dari fungsi terintegralkan-kuadrat 284O, 286U


----- rumus inversi untuk fungsi terdiferensialkan 283I; untuk fungsi bervariasi terbatas 283L


----- $\varhatf(y)=\lim_{\epsilon\downarrow 0}
  \bover1{\sqrt{2\pi}}\int_{-\infty}^{\infty}
  e^{-iyx}e^{-\epsilon x^2}f(x)dx$ a.e.\ 284M




\indexheader{bebas}








\indexheader{Freese}







\indexheader{Fubini} teorema Fubini \S252

----- $\int fd(\mu\times\nu)=\iint f(x,y)dxdy$ 252B

----- bila kedua faktor $\sigma$-hingga 252C, 252H

----- untuk fungsi indikator 252D, 252F



\indexheader{fungsi}


\indexheader{dasar} Teorema Dasar Kalkulus ($\bover{d}{dx}\int_a^xf=f(x)$ hampir di mana-mana) \S222; ($\int_a^bF'(x)dx=F(b)-F(a)$) 225E


\indexmedskip
















\indexmedskip




















\indexheader{Hahn} dekomposisi Hahn dari suatu fungsional aditif terhitung 231E





\indexheader{Hardy} Teorema Maksimal Hardy--Littlewood 286A








\indexheader{ukuran Hausdorff} ukuran Hausdorff (pada $\BbbR^r$) \S264

----- merupakan ukuran topologis 264E

----- ukuran Hausdorff berdimensi $r$ pada $\BbbR^r$ suatu kelipatan ukuran Lebesgue 264I












----- ukuran berdimensi $(r-1)$ pada $\BbbR^r$ 265F-265H








\indexmedskip

\indexiiheader{citra} ukuran citra 234C





\indexheader{tak tentu} integral tak tentu

----- bila diturunkan kembali menghasilkan fungsi asal 222E, 222H


----- untuk mengonstruksi ukuran 234I



\indexheader{bebas} variabel acak bebas \S272

----- distribusi bersama merupakan ukuran produk 272G


----- jumlah 272S, 272T, 272V, 272W

----- teorema limit \S273, \S274










regularitas dalam ukuran



----- (terhadap himpunan kompak) ukuran Lebesgue 134F; ukuran Radon \S256







\indexheader{integral}










\indexheader{integrasi} konstruksi integral fungsi bernilai real \S122

----- sebagai fungsional linear positif 122O

----- ----- bertindak pada $L^1(\mu)$ 242B

----- secara parsial 225F

----- karakterisasi fungsi terintegralkan 122P, 122R

----- pada subhimpunan 131D, 214E

----- fungsi dan integral dengan nilai dalam $[-\infty,\infty]$ 133A, 135F




























\indexmedskip

\indexiiheader{Jensen}


pertidaksamaan Jensen 233I-233J, 242K

----- dinyatakan dalam $L^1(\mu)$ 242K

\indexmedskip









\indexiiheader{Koml\'os} Teorema subbarisan Koml\'os 276H

\indexheader{Kwapien}


\indexmedskip

\indexheader{Teorema Kepadatan Lebesgue} Teorema Kepadatan Lebesgue (dalam $\Bbb R$) \S223

----- $\lim_{h\downarrow 0}\bover1{2h}\int_{x-h}^{x+h}f=f(x)$ a.e.\
223A


-----
$\lim_{h\downarrow 0}\bover1{2h}\int_{x-h}^{x+h}|f(x-y)|dy=0$ a.e.\ 223D


----- (dalam $\BbbR^r$) 261C, 261E





ukuran Lebesgue, konstruksi \S114, \S115

----- sifat-sifat lebih lanjut \S134



Teorema Konvergensi Terdominasi Lebesgue ($\int\lim=\lim\int$
untuk barisan fungsi yang didominasi) 123C

\indexheader{Levi} Teorema B. Levi ($\int\lim=\lim\int$ untuk barisan fungsi monoton) 123A; ($\int\sum=\sum\int$) 226E




\indexheader{pengangkatan}

















\indexheader{Lipschitz} fungsi Lipschitz \S262

----- terdiferensialkan hampir di mana-mana\ 262Q

\indexheader{dapat dilokalkan} ruang ukur yang dapat dilokalkan

----- merangkai fungsi terukur parsial 213N

----- yang tidak dapat dilokalkan secara ketat 216E















\indexheader{Loomis}


\indexheader{bawah}






\indexheader{Lusin}


teorema Lusin (fungsi terukur hampir kontinu)


----- (pada $\BbbR^r$) 256F





\indexmedskip














































\indexiiheader{martingal} martingal \S275

----- $L^1$-terbatas martingal konvergen hampir di mana-mana\ 275G


----- bila berbentuk $\Expn(X|\Sigma_n)$ 275H, 275I

\indexheader{terukur}




selubung terukur

----- sifat-sifat elementer 132E

----- eksistensi 213L

\wheader{terukur}{0}{0}{0}{36pt}

fungsi terukur

----- (bernilai real) \S121

----- ----- jumlah, produk, dan operasi lain pada berhingga banyak fungsi 121E

----- ----- limit, infimum, supremum 121F













\indexheader{ukuran}














ruang ukur \S112













\indexheader{monoton} Teorema Kelas Monoton (untuk aljabar himpunan) 136B

\indexheader{monoton} fungsi monoton

----- terdiferensialkan hampir di mana-mana\ 222A

\indexmedskip



\indexheader{sempit}












\indexheader{tak terukur} himpunan tak terukur (untuk ukuran Lebesgue) 134B

\indexheader{tak paradoksal}














\indexmedskip

\indexheader{keterurutan}














\indexiiheader{pengurutan} pengurutan ukuran 234P




\indexheader{luar} ukuran luar yang dikonstruksi dari ukuran 132A {\it dst.}

regularitas luar ukuran Lebesgue 134F

\indexmedskip

















\indexiiheader{Plancherel} teorema Plancherel (tentang deret dan transformasi Fourier dari fungsi terintegralkan-kuadrat) 282K, 284O










\indexheader{Poisson}




teorema Poisson (cacah kejadian langka yang bebas mempunyai distribusi yang mendekati Poisson) 285Q














\indexiiheader{produk} produk dua ruang ukur \S251

----- sifat dasar ukuran produk c.l.d.\ 251I

----- ukuran Lebesgue pada $\BbbR^{r+s}$ sebagai ukuran produk 251N


----- lebih dari dua ruang 251W

----- teorema Fubini 252B


----- teorema Tonelli 252G


----- dan ruang $L^1$ \S253

----- pemetaan bilinear kontinu pada $L^1(\mu)\times L^1(\nu)$ 253F


----- ekspektasi bersyarat pada suatu faktor 253H







produk sebarang banyak ruang probabilitas \S254

----- sifat dasar produk yang dilengkapkan 254F

----- karakterisasi melalui fungsi \imp\ 254G

----- produk subruang 254L

----- produk dari produk 254N

----- penentuan oleh subproduk terhitung 254O

----- subproduk dan ekspektasi bersyarat 254R



















\indexmedskip





















\indexmedskip

\indexiiheader{Rademacher} teorema Rademacher (fungsi Lipschitz terdiferensialkan hampir di mana-mana) 262Q


\indexiiheader{Radon} ukuran Radon


----- pada $\BbbR^r$ \S256

----- sebagai pelengkapan ukuran Borel 256C


----- ukuran integral tak tentu 256E


----- ukuran citra 256G


----- ukuran produk 256K




















Teorema Radon--Nikod\'ym (fungsi himpunan aditif yang benar-benar kontinu mempunyai fungsi kepadatan) 232E


----- dalam istilah $L^1(\mu)$ 242I





























integral berulang \S252







































indexvheader{Rothberger}




\indexmedskip































































\indexheader{Stone}
















teorema Stone--Weierstrass \S281

----- untuk subruang Riesz dari $C_b(X)$ 281A

----- untuk subaljabar dari $C_b(X)$ 281E

----- untuk *-subaljabar dari $C_b(X;\Bbb C)$ 281G



\indexheader{secara ketat} ukuran yang dapat dilokalkan secara ketat

----- syarat cukup agar dapat dilokalkan secara ketat 213O



\indexheader{hukum kuat} hukum kuat bilangan besar ($\lim_{n\to\infty}\bover1{n+1}\sum_{i=0}^n(X_i-\Expn(X_i))=0$ hampir di mana-mana) \S273


----- bila $\sum_{n=0}^{\infty}\Bover1{(n+1)^2}\Var(X_n)<\infty$ 273D


----- bila $\sup_{n\in\Bbb N}\Expn(|X_n|^{1+\delta})<\infty$ 273H


----- untuk yang berdistribusi identik $X_n$ 273I

----- untuk barisan selisih martingal 276C, 276F

----- konvergensi rata-rata untuk $\|\,\|_1$, $\|\,\|_p$ 273N
















\indexheader{subruang} ukuran subruang

----- untuk subruang terukur \S131

----- untuk subruang sebarang \S214




\indexiiheader{permukaan} ukuran permukaan dalam $\BbbR^r$ \S265







\indexmedskip











\indexiiheader{tensor} produk tensor ruang $L^1$ \S253








\indexiiheader{Tonelli} teorema Tonelli ($f$ adalah terintegralkan jika $\iint|f(x,y)|dxdy<\infty$) 252G





















\indexmedskip










\indexiiheader{seragam}




himpunan terintegralkan secara seragam dalam $L^1$ \S246

----- kriteria untuk keterintegralan seragam 246G



----- dan konvergensi dalam ukuran 246J

----- dan kekompakan lemah 247C









\indexmedskip




\indexiiheader{Vitali} teorema Vitali (untuk penutup oleh interval dalam $\Bbb R$) 221A

----- (untuk penutup oleh bola dalam $\BbbR^r$) 261B





\indexmedskip

\indexiiheader{lemah} kekompakan lemah dalam $L^1(\mu)$ \S247












\wheader{secara lemah}{0}{0}{0}{36pt}








\indexiiheader{Weyl} Teorema Ekuidistribusi Weyl 281N




\indexmedskip




\indexiiheader{$\scriptstyle\pmb{C}$} versi c.l.d.\ 213D {\it dst.}





























$L^0(\mu)$ (ruang kelas ekuivalensi fungsi terukur) \S241

----- sebagai ruang Riesz 241E





















$L^1(\mu)$ (ruang kelas ekuivalensi fungsi terintegralkan) \S242


----- kelengkapan norma 242F

----- kepadatan fungsi sederhana 242M

----- (untuk ukuran Lebesgue) kepadatan fungsi kontinu dan fungsi tangga 242O
















$L^p(\mu)$ (ruang kelas ekuivalensi fungsi yang pangkat ke-$p$-nya terintegralkan, dengan $1<p<\infty$) \S244


----- merupakan kisi Banach 244G

----- mempunyai dual $L^q$, dengan $\bover1p+\bover1q=1$ 244K

----- dan ekspektasi bersyarat 244M









$L^{\infty}(\mu)$ (ruang kelas ekuivalensi fungsi terukur terbatas) \S243


----- dualitas dengan $L^1$ 243F-243G

----- kelengkapan norma 243E

----- kelengkapan urutan 243H



















































\indexmedskip \indexheader{$\sigma$} aljabar-$\sigma$ himpunan \S111

----- yang dibangkitkan oleh keluarga-keluarga yang diberikan 136B, 136G



ukuran $\sigma$-hingga 215B



\indexheader{$\tau$}



















  

\discrpage

\wheader{Indeks umum}{24}{12}{2}{100pt}


\gdef\newparagraph{}
\gdef\bottomparagraph{}
\def\chaptername{Indeks}
\def\sectionname{Indeks umum}

\noindent{\bf\sectionname}
   \smallskip
   \let\headlinesectionname=\sectionname

\indexmedskip


Rujukan yang dicetak dengan huruf {\bf tebal} mengacu pada definisi;
rujukan dengan huruf {\it miring} merupakan rujukan sepintas.
Definisi yang ditandai dengan $\pmb{>}$ adalah definisi yang pemakaian saya
menyimpang secara berbahaya dari pemakaian sebagian penulis lain.

\indexmedskip



\indexheader{Abel} teorema Abel 224Yi




\indexiiheader{mutlak} keterjumlahan mutlak {\it 226Ac}


\indexiiheader{kontinu mutlak} fungsional aditif kontinu mutlak {\bf 232Aa}, 232B, 232D, 232F-232I, 232Xa, 232Xb, 232Xd, 232Xf, 232Xh, {\it 232Yk}, {\it 234Kb}, {\it 256J}, 257Xf



fungsi kontinu mutlak {\bf 225B}, 225C-225G, 225K-225O, 225Xa-225Xf, 225Xl-225Xn, 225Ya, 225Yc, 225Yd, {\it 232Xb}, {\it 233Xc}, {\it 244Yi}, 252Ye, 256Xg, 262Bc, 262Xl, 263J, {\bf 264Yp}, 265Ya, {\it 274Xb}, 282R, 282Yf, {\it 283Ci}, {\it 283Ya}, {\it 283Yb}



























\indexiiheader{teradaptasi} martingal teradaptasi {\bf 275A}


----- waktu henti {\bf 275L}


\indexheader{aditif} fungsional aditif pada aljabar himpunan {\it lihat} aditif hingga ({\bf 136Xg}, {\bf 231A}), aditif terhitung ({\bf 231C})




















\indexheader{adjoint} operator adjoint 243Fc, 243 {\it catatan}










\indexheader{aljabar} aljabar {\it lihat} aljabar himpunan ({\bf 136E}), aljabar Banach ({\bf 2A4Jb})), aljabar bernorma ({\bf 2A4J})


aljabar himpunan 113Yi, {\bf 136E}, 136F, 136G, 136Xg, 136Xh, 136Xk, 136Ya, 136Yb, {\bf 231A}, 231B, 231Xa;  

  {\it lihat juga } aljabar-$\sigma$ ({\bf 111A})






\indexheader{hampir} fungsi hampir kontinu 256F







hampir setiap, hampir di mana-mana {\bf 112Dd}








hampir pasti {\bf 112De}

\indexheader{berselang-seling} fungsional berselang-seling {\bf 132Yf}
















\indexheader{analitik} fungsi analitik (kompleks) 133Xc



\indexheader{angelik} ruang topologis angelik {\bf 2A5J}


\indexheader{antirantai}




\indexheader{aperiodik}






\indexheader{Archimedean}


ruang Riesz Archimedean {\bf 241F}, 241Yb, 242Xc



\indexiiheader{Archimedes} Archimedes 265Xf


\indexiiheader{luas} luas {\it lihat} ukuran permukaan





\indexiiheader{aritmetika} rata-rata aritmetika 266A















\indexiiheader{asimtotik} kepadatan asimtotik {\bf 273G}, 273Xo








\indexheader{secara asimtotik} terekuidistribusi secara asimtotik {\it lihat} terekuidistribusi ({\bf 281Yi})




\indexheader{atom}


atom (dalam ruang ukur) {\bf 211I}, {\it 211Xb}, {\it 246G}


\indexiiheader{atomik}


{atomik {\it lihat}}

atomik murni ({\bf 211K})


\indexheader{tanpa atom}




ukuran (ruang ukur) tanpa atom {\bf 211J}, 211Md, {\it 211Q}, {\it 211Xb}, 211Yd, 211Ye, {\it 212G}, 213He, 214K, 214Q, 215D, 215E, 215Xe, 215Xf, {\it 216A}, 216Ya, 216Ye, 234Be, 234Nf, {\it 234Yi}, 251U, 251Wo, 251Xs, 251Xt, 251Yd, 252Yo, 252Yq, 252Ys, 254V, 254Yh, 256Xd, 264Yg






















\indexiiheader{automorfisme} automorfisme {\it lihat}









automorfisme ruang ukur









\indexheader{aksioma} aksioma {\it lihat} masalah Banach--Ulam, pilihan ({\bf 2A1J}), pilihan terhitung

































































\indexmedskip





































\indexheader{bola} bola (dalam $\BbbR^r$) 252Q, 252Xi, 264H;
  
  {\it lihat juga} sfera


\indexheader{Banach} aljabar Banach $\pmb{>}${\bf 2A4Jb}




kisi Banach {\bf 242G}, 242Xc, 242Yc, 242Ye, 243E, 243Xb; {\it lihat juga} $L^p$






ruang Banach {\it 231Yh}, {\it 262Ya}, {\bf 2A4D}, 2A4E, 2A4I; {\it lihat juga} aljabar Banach ({\bf 2A4Jb}), kisi Banach ({\bf 242G}), ruang Banach separabel


masalah Banach--Ulam 232Hc












































\indexiiheader{Berry-Ess\'een} teorema Berry--Ess\'een 274Hc, 285 {\it catatan}








\indexheader{bidual}




\indexheader{Bienaym\'e} kesamaan Bienaym\'e 272S


\indexheader{bilateral}


\indexheader{bilinear} pemetaan bilinear \S253 ({\bf 253A}), {\it 255 catatan}












\indexheader{Bochner} integral Bochner {\bf 253Yf}, 253Yg, 253Yi


teorema Bochner 285Xu




\indexheader{Boolean}



























\indexheader{Borel} aljabar Borel {\it lihat} aljabar-$\sigma$ Borel ({\bf 111Gd}, {\bf 135C}, {\bf 256Yf}) 





lema Borel--Cantelli 273K


















fungsi terukur Borel {\bf 121C}, 121D, 121Eg, 121H, 121K, {\bf 121Yd}, 134Fd, 134Xg, 134Yt, {\bf 135Ef}, 135Xc, 135Xe, 225H, 225J, {\it 225Yg}, 233Hc, 241Be, 241I, 241Xd, 256M, 262Yd

ukuran Borel (pada $\Bbb R$) 211P, 216A 






himpunan Borel dalam $\Bbb R$, $\BbbR^r$ {\bf 111G}, {\it 111Yd}, 114G, 114Yd, 115G, 115Yb, {\it 115Yd}, {\it 121Ef}, {\it 121K}, 134F, 134Xd, {\it 135C}, 136D, 136Xj, 254Xi, 225J, 264E, 264F, 264Xb, 266Bc
  






aljabar-$\sigma$ Borel (atas subhimpunan $\Bbb R^r$) {\bf 111Gd}, {\it 114Yg-114Yi}, 121J, 121Xd, 121Xe, 121Yc, 121Yd, {\it 212Xc}, {\it 212Xd}, {\it 216A}, 251M;

----- (dari ruang lain) {\bf 135C}, 135Xb, {\it 254Xi}, {\bf 256Yf}, 271Ya






\indexheader{terbatas} operator bilinear terbatas {\bf 253Ab}, 253E, 253F, 253L, 253Xb, 253Yf


operator linear terbatas 242Je, 253Ab, 253F, 253Gc, 253L, 253Xc, 253Yf, 253Yj, {\bf 2A4F}, 2A4G-2A4I, 2A5If; {\it lihat juga} $\eurm B(U;V)$ ({\bf 2A4F})



himpunan terbatas (dalam $\BbbR^r$) {\bf 134E}; (dalam ruang bernorma) {\bf 2A4Bc}; (dalam ruang topologis linear) 245Yf; {\it lihat juga} terbatas menurut urutan ({\bf 2A1Ab})
  

fungsi berdukungan terbatas {\it lihat} dukungan kompak


fungsi bervariasi terbatas \S224 ({\bf 224A}, {\bf 224K}), 225Cb, 225M, {\it 225Oc}, {\it 225Xf}, {\it 225Xm}, 225Yc, 226Bc, 226C, 226Yd, 263Ye, 264Yp, 265Yb, 282M, 282O, 283L, 283Xj, 283Xk, {\it 283Xm}, {\it 283Xn}, {\it 283Ya}, {\it 283Yb}, 284Xl, {\it 284Yd}


\indexheader{pembatasan}

















\indexheader{Breiman}



















\indexiiheader{Brunn} pertidaksamaan Brunn--Minkowski 266C













\indexmedskip















\indexheader{Cantor} fungsi Cantor {\bf 134H}, 134I, 222Xd, 225N, 226Cc, 262K, 264Xf, 264Yn

ukuran Cantor {\bf 256Hc}, 256Ia, 256Xk, 256Yd, 264Ym


himpunan Cantor {\bf 134G}, 134H, 134I, 134Xf, 256Hc, 256Xk, 264J, 264Ym, 264Yn; {\it lihat juga} $\{0,1\}^I$









\indexheader{\Caratheodory} Ruang ukur lengkap \Caratheodory\ {\it lihat} lengkap ({\bf 211A})


Metode \Caratheodory: konstruksi ukuran 113C, 113D, 113Xa, 113Xd, 113Xg, 113Yc, 113Yk, 114E, 114Xa, 115E, {\it 121Ye}, 132Xc, 136Ya, 212A, {\it 212Xf}, 213C, 213Xa, {\it 213Xb}, {\it 213Xd}, 213Xf, {\it 213Xg}, 213Xi, 213Yb, 214H, 214Xb, 216Xb, {\it 251C}, {\it 251Wa}, {\it 251Xe}, {\it 264C}, {\it 264K}


\indexheader{kardinal} kardinal {\bf 2A1Kb}, 2A1L





















\indexheader{Carleson} teorema Carleson 282K {\it komentar}, 282 {\it catatan}, {\it 284Yg}, 286U, 286V











\indexheader{kategori}


\indexheader{Cauchy} distribusi Cauchy {\bf 285Xp}


filter Cauchy {\bf 2A5F}, 2A5G


pertidaksamaan Cauchy 244Eb




barisan Cauchy {\it 242Yc}, {\bf 2A4D}









\indexheader{selularitas}






















\indexheader{pusat} Teorema Limit Pusat 274G, 274I-274K, 274Xd-274Xh, {\it 274Xj}, {\it 274Xk}, 285N, 285Xq, 285Yo


\indexiiheader{Ces\`aro} jumlah Ces\`aro 273Ca, 282Ad, 282N, 282Xn







\indexheader{rantai}




aturan rantai untuk turunan Radon--Nikod\'ym 235Xh







\indexiiheader{penggantian} penggantian variabel dalam pengintegralan \S235, 263A, 263D, 263F, 263G, 263I, 263J, 263Xc, 263Xe, 263Yc





\indexheader{karakteristik} fungsi karakteristik (dari suatu himpunan) {\it lihat} fungsi indikator ({\bf 122Aa})

----- (dari suatu distribusi probabilitas) \S285 ({\bf 285Aa})


----- (dari suatu variabel acak) \S285 ({\bf 285Ab})


\indexheader{muatan}


\indexheader{dapat diberi muatan}


\indexheader{Chebyshev} pertidaksamaan Chebyshev 271Xa


\indexheader{pilihan} aksioma pilihan 134C, 1A1G, 254Da, {\bf 2A1J}, 2A1K-2A1M; {\it lihat juga} pilihan terhitung

fungsi pilihan {\bf 2A1J}















\indexheader{lingkaran} grup lingkaran 255M, {\it 255Ym}




\indexheader{buka-tutup}



\indexheader{tertutup}


selubung konveks tertutup 253Gb, {\bf 2A5Eb}




interval tertutup (dalam $\Bbb R$ atau $\BbbR^r$) 114G, 115G, 1A1A; (dalam ruang terurut parsial umum) {\bf 2A1Ab}
  

himpunan tertutup (dalam ruang topologis) 134Fb, 134Xd, {\bf 1A2E}, 1A2F, 1A2G, 2A2A, {\bf 2A3A}, 2A3D, {\it 2A3Nb}
  





\indexheader{penutupan} penutupan suatu himpunan {\bf 2A2A}, 2A2B, {\bf 2A3D}, 2A3Kb, 2A5E; {\it lihat juga} penutupan esensial ({\bf 266B})


\indexheader{club (tertutup tak terbatas)}


\indexheader{gugus} titik gugus (dari suatu barisan) {\bf 2A3O}






\indexheader{koarea} Teorema Koarea {\it 265 catatan}





























































\indexheader{kohingga} kohingga {\it lihat} aljabar hingga-kohingga ({\bf 231Xa})






















\indexheader{komutatif} komutatif (gelanggang, aljabar) {\bf 2A4Jb}


aljabar Banach komutatif 224Yb, 243Dd, 255Xi, 257Ya




\indexheader{komutator}




\indexheader{saling komutatif}


\indexheader{kompak}














himpunan kompak (dalam ruang topologis) 247Xc, {\bf 2A2D}, 2A2E-2A2G, $\pmb{>}${\bf 2A3N}, 2A3R;  
  
  {\it lihat juga} relatif kompak ({\bf 2A3Na}), relatif kompak lemah ({\bf 2A5Id}), kompak lemah ({\bf 2A5Ic})




fungsi dengan dukungan kompak 242Xh, 256Be, {\it 256D}, 262Ye-262Yh; {\it lihat juga} $C_k(X)$, $C_k(X;\Bbb C)$


ukuran dengan dukungan kompak 284Yi





ruang topologis kompak {\bf $\pmb{>}$2A3N}



\indexheader{kompaktifikasi}














\indexheader{komplemen}


\indexheader{komplementer}


\indexheader{terkomplemen}





\indexheader{lengkap}




ruang topologis linear lengkap 245Ec, {\bf 2A5F}, 2A5H


ukuran (ruang ukur) lengkap yang ditentukan secara lokal 213C, 213D, 213H, 213J, 213O, 213Xg, 213Xk, 213Xl, 213Ye, 213Yf, 214I, 216D, 216E, 234Nb, 251Yc, 252B, 252D, 252N, 252Xc, 252Ye, 252Ym, 252Yq-252Ys, 252Yv, 253Yj, 253Yk; {\it lihat juga} versi c.l.d.\ ({\bf 213E})





ukuran (ruang ukur) lengkap {\bf 112Df}, 113Xa, 122Ya, {\bf 211A}, 211M, 211N, 211R, {\it 211Xc}, {\it 211Xd}, \S212, 214I, {\it 214K}, {\it 216A}, {\it 216C-216E}, {\it 216Ya}, 216Ye, 234Ha, {\it 234I}, 234Ld, 234Ye, 235Xl, {\it 254Fd}, 254G, {\it 254J}, 256Ad, 257Yb, {\it 264Dc}; {\it lihat juga} ukuran lengkap yang ditentukan secara lokal
  


ruang metrik lengkap {\it 224Ye}; {\it lihat juga} ruang Banach ({\bf 2A4D})


ruang bernorma lengkap {\it lihat} ruang Banach ({\bf 2A4D}), kisi Banach ({\bf 242G})


ruang Riesz lengkap {\it lihat} Dedekind lengkap ({\bf 241Fc})










\indexheader{dilengkapkan} ukuran integral tak tentu yang dilengkapkan {\it lihat} ukuran integral tak tentu ({\bf 234J})


\indexheader{sepenuhnya}
















\indexheader{pelengkapan}






{pelengkapan } (dari suatu ukuran (ruang ukur)) \S212 ({\bf 212C}), 213Fa, 213Xa, 213Xb, 213Xi, 214Ib, 214Xi, 216Yd, {\it 232Xe}, 234Ba, 234Ke, {\it 234Lb}, 234Xc, 234Xl, 234Ye, 234Yo, {\it 235D}, 241Xb, 242Xb, 243Xa, 244Xa, 245Xb, 251T, 251Wn, 251Xr, 252Ya, 254I, 256C




















\indexheader{kompleksifikasi}


\indexheader{bernilai kompleks} fungsi bernilai kompleks \S133

\indexheader{komponen} komponen (dalam ruang topologis) {\it 111Ye}

























\indexheader{bersyarat}




ekspektasi bersyarat {\bf 233D}, 233E, 233J, 233K, 233Xg, 233Yc, 233Ye, 235Yb, 242J, 246Ea, 253H, 253Le, 275Ba, 275H, 275I, 275K, 275Ne, 275Xj, 275Yb, 275Yo, {\it 275Yp}


operator ekspektasi bersyarat {\bf 242Jf}, 242K, 242L, 242Xe, 242Xj, 242Yk, 243J, 244M, 244Yl, 246D, 254R, 254Xq, 275Xd, 275Xe










\indexheader{konegligibel} himpunan konegligibel {\bf 112Dc}, 211Xg, 212Xi, 214Cc, 234Hb

\indexheader{terhubung} himpunan terhubung {\bf 222Yb}







\indexheader{konjugat}










\indexheader{konsisten}














\indexheader{kontinuitas} titik-titik kontinuitas {\it 224H}, {\it 224Ye}, 225J


\indexheader{kontinu} fungsi kontinu 121D, 121Yd, {\it 224Xl}, {\it 262Ia}, {\bf 2A2C}, 2A2G, {\bf 2A3B}, 2A3H, 2A3Nb, 2A3Qb





fungsional linear kontinu {\it 284Yj}; {\it lihat juga} ruang linear dual ({\bf 2A4H})


operator linear kontinu 2A4Fc; {\it lihat juga} operator linear terbatas ({\bf 2A4F})


ukuran kontinu {\it lihat} tanpa atom ({\bf 211J}), kontinu mutlak ({\bf 232Aa})


kontinu dari kanan {\bf 224Xm}





kontinu {\it lihat juga }

semi-{\vthsp}kontinu ({\bf 225H})



\indexheader{kontinuum} kontinuum {\it lihat} $\frak c$ ({\bf 2A1L})









\indexheader{konvergensi} konvergensi dalam rata-rata (dalam $\eusm L^1(\mu)$ atau $L^1(\mu)$) {\bf 245Ib}


\indexiiheader{konvergensi dalam ukuran} konvergensi dalam ukuran (dalam $L^0(\mu)$) \S245 ($\pmb{>}${\bf 245A}), 246J, 246Yc, 247Ya, 253Xe, 255Yh, 285Yk


----- ----- (dalam $\eusm L^0(\mu)$) \S245 ($\pmb{>}${\bf 245A}), 271Yd, 272Yd, 274Yf, 285Yt


----- ----- (dalam aljabar himpunan terukur) {\bf 232Ya}




----- ----- (untuk barisan) 245Ad, 245C, 245H-245L, 245Xd, 245Xf, 245Xl, 245Yh, 246J, 246Xh, 246Xi, 271L, 273Ba, 275Xl, 276Yf




\indexheader{konvergen} konvergen hampir di mana-mana 245C, 245K, 273Ba, 276G, 276H; {\it lihat juga} hukum kuat


filter konvergen {\bf 2A3Q}, 2A3S, 2A5Ib;  




{----- } barisan {\bf 135D}, 245Yi, {\bf 2A3M}, 2A3Sg; {\it lihat juga} konvergensi dalam ukuran ({\bf 245Ad})
  



\indexheader{konveks} fungsi konveks {\bf 233G}, 233H-233J, 233Xb-233Xf ({\bf 233Xd}), 233Xh, 233Ya, 233Yc, 233Yd, 233Yg, 233Yi, 233Yj, 242K, 242Yi, 242Yj, 242Yk, {\it 244Xm}, {\it 244Yh}, 244Ym, {\it 255Yj}, {\it 275Yh}; {\it lihat juga} konveks-tengah ({\bf 233Ya})


selubung konveks {\bf 2A5E}; {\it lihat juga} selubung konveks tertutup ({\bf 2A5Eb})






himpunan konveks {\bf 233Xd}, {\bf 244Yk}, 233Yf-233Yj, {\it 262Xh}, 264Xg, 266Xc, {\bf 2A5E}











\indexheader{konvolusi} konvolusi dalam $L^0$ {\bf 255Xh}, 255Xi, 255Xj, 255Yh, 255Yj


konvolusi fungsi {\bf 255E}, 255F-255K, {\bf 255O}, 255Xa-255Xg, 255Ya, 255Yc, 255Yd, 255Yf, 255Yg, 255Yk, 255Yl, {\bf 255Yn}, 262Xj, 262Ye, 262Yf, 262Yi, 263Ya, 282Q, 282Xt, 283M, 283Xl, 283Wg, 283Wh, 283Wj, 284J, 284K, 284O, 284Wf, 284Wi, 284Xb, 284Xe


konvolusi ukuran \S257 ({\bf 257A}), 272T, 285R, 285Yp


konvolusi ukuran dan fungsi {\it 257Xe}, 284Xp, 284Yi


konvolusi barisan 255Xk, {\it 255Yo}, 282Xq








\indexheader{terhitung} himpunan terhitung {\bf 111F}, 114G, 115G, \S1A1, 226Ya







aksioma pilihan terhitung {\it 134C}, 211P






aljabar terhitung-koterhitung {\bf 211R}, 211Ya, {\it 232Hb}


ukuran terhitung-koterhitung {\bf 211R}, {\it 232Hb}, 252L






sifat supremum terhitung (dalam suatu ruang Riesz) {\bf 241Ye}, 242Yd, 242Ye, 244Yb




\indexiiheader{secara terhitung} fungsional aditif terhitung (pada suatu aljabar-$\sigma$ himpunan) {\bf 231C}, 231D-231F, 231Xf, {\it 231Xg}, 231Ya, 231Yb, 231Yf-231Yh, \S232, 246Yg, 246Yi





































\indexheader{pencacahan} ukuran cacah {\bf 112Bd}, {\it 122Xd}, 122 {\it catatan}, 211N, 211Xa, {\it 213Xb}, 226Ab, 241Xa, 242Xa, 243Xl, 244Xh, 244Xn, 245Xa, 246Xc, 246Xd, 251Xc, 251Xi, {\it 252K}, 255Yo, {\it 264Db}
  










\indexheader{penutup} penutup {\it lihat} selubung terukur ({\bf 132D})


Lema Penutup {\it lihat} Lema Penutup Lebesgue


















teorema penutup 221A, 261B, 261F, {\it 261Xc}, 261Ya, 261Yi, 261Yk




















\indexiiheader{silinder} silinder 265Xf; (dalam `silinder terukur') {\bf 254Aa}, 254F, {\it 254G}, {\it 254Q}, 254Xa






\indexmedskip








\indexheader{De Morgan}




\indexheader{desimal}
pengembangan desimal 273Xg


\indexheader{dapat didekomposisi}


ukuran (ruang ukur) yang dapat didekomposisi {\it lihat} dapat dilokalkan secara ketat ({\bf 211E})


\indexheader{dekomposisi}
dekomposisi (ruang ukur) {\bf 211E}, {\it 211Yd}, 213Ob,
{\it 213Xj}, 214Ia, 214L, 214N, {\it 214Xh} 


-----  {\it lihat juga }
  dekomposisi Hahn (231F), dekomposisi Jordan (231F), dekomposisi Lebesgue (232I, {\bf 226C})


\indexheader{menurun}
susunan ulang menurun (dari suatu fungsi) {\bf 252Yo}




\indexheader{Dedekind}


{himpunan terurut lengkap Dedekind } {\it 135Ba}, 234Yo 


----- ----- ruang Riesz {\bf 241Fc}, 241G, 241Xf, 242H, 242Yc,
243Hb, 243Xj, 244L






\wheader{Dedekind}{0}{0}{0}{72pt}





ruang Riesz lengkap-$\sigma$ Dedekind {\bf 241Fb}, 241G, 241Xe, 241Yb, 241Yh, 242Yg, 243Ha, 243Xb




\indexiiheader{delta}
fungsi delta {\it lihat} fungsi delta Dirac ({\bf 284R})


sistem delta {\it lihat} sistem $\Delta$




\indexiiheader{Denjoy}
Teorema Denjoy-Young-Saks 222L


\indexiiheader{padat}
himpunan padat dalam ruang topologis 136H,
242Mb, 242Ob, 242Pd, 242Xi, 243Ib, 244H, 244Pb, 244Xk, 244Yj, 254Xp, 281Yc,
{\bf 2A3U}, 2A4I






\indexiiheader{kerapatan}




fungsi kerapatan (variabel acak) {\bf 271H},
271I-271K, 
271Xc-271Xe, 
272U, 272Xd, 272Xj; 
  {\it lihat juga} fungsi kerapatan normal ({\bf 274A}), turunan Radon-Nikod\'ym ({\bf 232Hf}, {\bf 234J})


titik kerapatan {\bf 223B}, 223Xc, {\it 223Yb}, 266Xb


topologi kerapatan {\bf 223Yb}, 223Yc, 223Yd, {\bf 261Yf},
261Yl


kerapatan {\it lihat juga}
  kerapatan asimtotik ({\bf 273G}), Teorema Kerapatan Lebesgue (223A)









\indexiiheader{turunan}
turunan {\it lihat} turunan Dini ({\bf 222J})


\indexheader{turunan}
turunan fungsi (satu variabel) 222C,
222E-222J, 
222Yd, 225J, {\it 225L}, {\it 225Of}, {\it 225Xb}, {\it 226Be}, 282R; (banyak variabel) {\bf 262F}, 262G, 262P, 263Ye; 
  {\it lihat} turunan parsial












\indexheader{determinan}
determinan matriks 2A6A


\indexheader{ditentukan}
ditentukan {\it lihat} ruang ukur yang ditentukan secara lokal ({\bf 211H}),
himpunan terabaikan yang ditentukan secara lokal ({\bf 213I})


ditentukan oleh koordinat (dalam `$W$ ditentukan oleh koordinat dalam
$J$') {\bf 254M}, 254O, 254R-254T, 
254Xq, 254Xs





\indexheader{tangga Setan}
Tangga Setan {\it lihat} fungsi Cantor ({\bf 134H})











\indexheader{keterdiferensialan}
titik-titik keterdiferensialan 222H, 225J


\indexheader{dapat didiferensialkan}
fungsi yang dapat didiferensialkan (satu variabel) 123D,
  {\it 222A}, {\it 224I}, {\it 224Kg}, {\it 224Yc}, {\it 225L}, {\it
225Of}, {\it 225Xb}, {\it 225Xm}, 233Xc, 252Ye, 255Xd, 255Xe,
265Xd, {\it 274E}, 282L, 282Rb, 282Xs,
283I-283K, 
{\it 283Xm}, 284Xc, 284Xl; (banyak variabel) {\bf 262Fa}, 262Gb, 262I, 262Xg,
262Xi-262Xk, 
263Yf, ;  
  {\it lihat juga} turunan


`terdiferensialkan relatif terhadap domainnya' 222L, {\bf 262Fb}, 262I,
262M-262Q, 
262Xd-262Xf, 
262Yc, 263D, 263I, 263Xc, 263Xd, 263Yc, 265E, 282Xk


\indexheader{mendiferensialkan}
diferensiasi di bawah tanda integral 123D

\indexheader{tersebar}
ukuran terdifusi {\it lihat} ukuran tanpa atom ({\bf 211J})


\indexheader{dilatasi}
dilatasi 284Xd, \S284 {\it catatan}, {\bf 286C}


\indexheader{dimensi}
dimensi {\it lihat} dimensi Hausdorff ({\bf 264 catatan})


\indexiiheader{Dini}
turunan Dini {\bf 222J}, 222L, 222Ye, 225Yg




\indexheader{Dirac}
ukuran Dirac {\bf 112Bd},
  257Xa, 274Lb, 284R, 284Xo, 284Xp, 285H, 285Xs

\indexheader{langsung}
citra langsung (dari suatu himpunan oleh fungsi atau relasi) 1A1B

jumlah langsung ruang ukur {\bf 214L}, 214M,
214Xh-214Xk, 
241Xg, 242Xd, 243Xe, {\it 244Xf}, {\it 245Yh}, 251Xi, 251Xj






\indexiiheader{terarah}
himpunan terarah
{\it lihat} terarah ke bawah ({\bf 2A1Ab}),
terarah ke atas ({\bf 2A1Ab})


\indexiiheader{Dirichlet}
kernel Dirichlet {\bf 282D};
{\it lihat juga} kernel Dirichlet termodifikasi ({\bf 282Xc})















\indexheader{disintegrasi}
disintegrasi suatu ukuran 123Ye


\indexheader{saling lepas}
keluarga himpunan yang saling lepas {\bf 112Bb}



teorema barisan saling lepas (dalam ruang Riesz topologis) 246G,
{\it 246Ha}, 246Yd-246Yf, 
246Yj






\indexheader{distribusi}
distribusi {\it lihat} distribusi Schwartz, distribusi tempered


distribusi variabel acak 241Xc,
271E-271G, 
272G, 272T, 272Xe, 272Ya, 272Yb, 272Yf, 285H, 285Xj;  
   
  {\it lihat juga} distribusi Cauchy ({\bf 285Xp}), distribusi empiris ({\bf 273 catatan}), distribusi Poisson ({\bf 285Xr})


----- dari suatu keluarga berhingga variabel acak  271B, {\bf 271C},
271D-271I, 
272G, 272Yb, 272Yc, 285Ab, 285C, 285Mb




fungsi distribusi variabel acak {\bf $\pmb{>}$271G}, 271L,
271Xb, 271Yb-271Yd, 
272Xe, 272Ya, 273Xh, 273Xi,
274F-274L, 
274Xd, 274Xe, 274Xh, 274Xi, 274Yc, 274Ye, 285P; 
  {\it lihat juga} fungsi distribusi normal ({\bf 274Aa})


\indexheader{distributif}








\indexheader{divergensi}












\indexheader{terdominasi}
Teorema Konvergensi Terdominasi {\it lihat} Teorema Konvergensi Terdominasi Lebesgue ({\bf 123C})




\indexheader{Doob}
Teorema Konvergensi Martingal Doob {\bf 275G}


ketaksamaan maksimal Doob 275D, 275Xb


\indexheader{ganda}


\indexheader{berganda}






















himpunan terurut parsial yang terarah ke bawah
{\bf 2A1Ab}














\indexheader{dual}







ruang bernorma dual 243G, 244K, {\bf 2A4H}












\indexheader{Dunford}




\indexheader{diadik}








\indexheader{Dye}





\indexheader{Dynkin}
kelas Dynkin {\bf 136A}, 136B, 136Xb




\indexmedskip

\indexheader{Eberlein}


Teorema Eberlein {\it 2A5J}















\indexheader{Egorov}
Teorema Egorov 131Ya, 215Yb
  





















\indexiiheader{empiris}
distribusi empiris 273Xi, {\bf 273 catatan}















\indexheader{entropi}




\indexheader{mengenumerasi}


\indexheader{selubung}
selubung {\it lihat}
selubung terukur ({\bf 132D})


















\indexiiheader{terekuidistribusi}
barisan terekuidistribusi (dalam ruang probabilitas topologis) 281N,
{\bf 281Yi}, {\it 281Yj}, 281Yk; 
  {\it lihat juga}
  terdistribusi baik ({\bf 281Ym})

  

\indexiiheader{ekuidistribusi}
Teorema Ekuidistribusi {\it lihat} Teorema Ekuidistribusi Weyl (281N)







\indexiiheader{ekuipoten}
himpunan ekuipoten 2A1G, 2A1Ha; 
  {\it lihat juga} kardinal ({\bf 2A1Kb}










\indexheader{ekuivalen}
ekuivalen {\bf 121B}, {\bf 212B}
























\indexiiheader{esensial}


penutupan esensial {\it 261Yj}, {\bf 266B}, 266Xa, 266Xb




supremum esensial keluarga himpunan terukur {\bf 211G}, 213K,
215B, 215Cb


----- dari suatu fungsi bernilai real {\bf 243D}, 243I, 255K


\indexheader{secara esensial}
fungsi berbatas esensial {\bf 243A}






\indexiiheader{Etemadi}
Lema Etemadi 272V




\indexheader{Euklides}
metrik Euklides (pada $\BbbR^r$) {\bf 2A3Fb}


topologi Euklides \S1A2, \S2A2, {\bf 2A3Ff}, 2A3Tc



\indexheader{genap}
fungsi genap 255Xb, 283Yc, 283Yd




\indexheader{dapat dipertukarkan}
keluarga variabel acak yang dapat dipertukarkan {\bf 276Xg}


\indexheader{pertukaran}


\indexheader{penghabisan}
prinsip penghabisan 215A, 215C, 215Xa, 215Xb,
{\it 232E}











\indexiiheader{ekspektasi}
ekspektasi (suatu distribusi) {\bf 271F}


----- (suatu variabel acak) {\bf 271Ab}, 271E, 271F, 271I,
{\it 271Xa}, 271Ye,
272R, 272Xb, {\it 272Xi}, 274Xb, 285Ga, {\it 285Xr}, 285Xh;
  {\it lihat juga} ekspektasi bersyarat ({\bf 233D})





\indexheader{secara eksponensial}


\indexheader{eksponensiasi}


\indexheader{diperluas}


garis real diperluas {\it 121C}, \S135

\indexheader{perluasan}




perluasan fungsional aditif berhingga 113Yi



perluasan ukuran 113Yc, 113Yh, 132Yd,
  212Xh, 214P, 214Xm, 214Xn,
214Ya-214Yc;  
  {\it lihat juga} pelengkapan ({\bf 212C}), versi c.l.d.\ ({\bf 213E})
  

















\indexmedskip




\indexiiheader{adil}
probabilitas koin adil 254J










\indexheader{Fatou}
Lema Fatou {\bf 123B}, 133Kb, 135Gb, 135Hb

norma Fatou pada ruang Riesz 244Yg








\indexheader{Federer}


\indexheader{Fej\'er}
integral Fej\'er {\bf 283Xf}, 283Xh-283Xj 


kernel Fej\'er {\bf 282D}


jumlah Fej\'er {\bf 282Ad}, 282B-282D, 
282G-282I, 
282Yc





\indexheader{Feller}
Feller, W.\  bab\ 27 {\it pengantar}













\indexheader{medan}
medan (himpunan) {\it lihat} aljabar ({\bf 136E})

\indexheader{filter}
filter 211Xg, {\bf 2A1I}, 2A1N, 2A1O;  
  {\it lihat juga}
filter Cauchy ({\bf 2A5F}), filter konvergen ({\bf 2A3Q}), filter Fr\'echet ({\bf 2A3Sg}), ultrafilter ({\bf 2A1N})









\indexiiheader{filtrasi}
filtrasi (dari aljabar-$\sigma$) {\bf 275A}




\indexheader{halus}







\indexheader{berhingga}
aljabar berhingga-koberhingga {\bf 231Xa}, 231Xc


















\indexheader{aditif berhingga}
fungsi (atau fungsional) aditif berhingga pada suatu aljabar himpunan 136Xg, 136Ya,
136Yb,
  {\bf 231A}, 231B, 231Xb-231Xf, 
231Ya-231Yh, 
232A, 232B, 232E, 232G, 232Ya, 232Ye, 232Yg, 243Xk; 
  {\it lihat juga} aditif terhitung
  



\vindexheader{aditif berhingga}{36}

























































\indexheader{Fourier}
koefisien Fourier {\bf 282Aa}, 282B, 282Cb, 282F, 282Ic, 282J,
282M, 282Q, 282R, 282Xa, 282Xg, 282Xq, 282Xt, 282Ya, 283Xt, 284Ya, 284Yg


rumus integral Fourier {\bf 283Xm}


deret Fourier {\it 121G}, \S282 ({\bf 282Ac})

jumlah Fourier {\bf 282Ab}, 282B-282D, 
282J, 282L, 282O, 282P, 282Xi-282Xk, 
282Xp, 282Xt, 282Yd, 286V, 286Xb


transformasi Fourier {\bf 133Xd}, {\bf 133Yc},
  \S283 ({\bf 283A}, {\bf 283Wa}), \S284 ({\bf 284H}, {\bf 284Wd}),
285Ba, 285D, 285Xe, 285Ya, 286U, 286Ya


transformasi Fourier-Stieltjes 
  {\it lihat} fungsi karakteristik ({\bf 285A})












\indexheader{Fr\'echet}
filter Fr\'echet {\bf 2A3Sg}






\indexheader{bebas}



















\indexheader{Fremlin}




\indexheader{Frol\'\i k}


\indexheader{Frostman}




\indexiiheader{Fubini}


Teorema Fubini 252B, 252C, 252H, 252R, 252Yb, 252Yc




\indexheader{penuh}




ukuran luar penuh {\bf 132F}, 132Yd, 134D, 134Yt,
  212Eb, 214F, 214J, 214Yd, 234F, 234Xa, 241Yg, 243Ya, 254Yh














\indexheader{fungsi}
fungsi 1A1B



\indexheader{fundamental}
Teorema Dasar Kalkulus 222E, 222H, {\it 222I}, 225E


Teorema Dasar Statistika 273Xi, 273 {\it catatan}



\indexmedskip




\indexheader{Gaifman}








\indexheader{permainan}


\indexheader{gamma}


fungsi gamma {\bf 225Xh}, {\it 225Xi}, 252Xi, 252Yf, 252Yu,
255Xg









\indexheader{tolok}




\indexheader{Gauss}




\indexheader{Gaussian}
distribusi Gauss
  {\it lihat} distribusi normal baku ({\bf 274Aa})








variabel acak Gauss {\it lihat} variabel acak normal ({\bf 274Ad})





\indexheader{dibangkitkan}






aljabar himpunan yang dibangkitkan (atau aljabar-$\sigma$ yang \nobreak dibangkitkan) {\bf 111Gb}, 111Xe, 111Xf,
{\it 121J}, 121Xd, 121Yc, 136B, 136C, 136G, 136H, 136Xc, 136Xk,
136Yb, 136Yc,
  251D, 251Xa, 272C



\indexheader{pembangkit}




\indexiiheader{geometrik}
rata-rata geometrik 266A


\indexheader{Geschke}














Teorema Glivenko-Cantelli {\it 273 catatan}










\indexheader{G\"odel}







\indexheader{gradien}


\indexheader{graf}
grafik fungsi 264Xf, 265Yb












\indexheader{grup}
grup 255Yn, 255Yo;
  {\it lihat juga} grup lingkaran












\indexmedskip


















\indexiiheader{Hahn}




dekomposisi Hahn (dari fungsional aditif terhitung)
231Eb, 231F


----- {\it lihat juga} Teorema Vitali-Hahn-Saks (246Yi)














\indexheader{setengah}
interval setengah terbuka (di $\Bbb R$ atau $\BbbR^r$) {\bf 114Aa}, 114G,
{\it 114Yj}, {\bf 115Ab}, 115Xa, {\it 115Xc}, {\it 115Yd}

setengah-ruang (dalam $\BbbR^r$) 285Xo


















\indexiiheader{Hanner}
ketaksamaan Hanner {\bf 244O}, 244Ym




\indexheader{Hardy}
Teorema Maksimal Hardy-Littlewood 286A












\indexheader{Hausdorff}


dimensi Hausdorff {\bf 264 catatan}






ukuran Hausdorff \S264
({\bf 264C}, {\bf 264Db}, {\bf 264K}, {\bf 264Yo}),
265Xd, 265Yb; 
  {\it lihat juga} ukuran Hausdorff ternormalisasi ({\bf 265A})


metrik Hausdorff (pada ruang himpunan tertutup) {\bf 246Yb}


ukuran luar Hausdorff \S264 ({\bf 264A}, {\bf 264K}, {\bf 264Yo})






topologi Hausdorff {\bf 2A3E}, 2A3L, 2A3Mb, 2A3S






























\indexheader{secara herediter}

















\indexheader{Hilbert}
ruang Hilbert 244Na, 244Yk
















\indexiiheader{Hoeffding}
Teorema Hoeffding 272W, 272Xk, 272Xl


\indexiiheader{H\"older}
fungsi kontinu H\"older {\bf 282Xj}, 282Yb


ketaksamaan H\"older {\bf 244Eb}, 244Yc, 244Ym














\indexheader{homomorfik}


\indexheader{homomorfisme}








\indexheader{selubung}
selubung {\it lihat} selubung konveks ({\bf 2A5Ea}), selubung konveks tertutup ({\bf 2A5Eb})


\indexheader{Humpty}















\indexmedskip

\indexheader{ideal}
ideal dalam suatu aljabar himpunan {\bf 214O}, {\bf 232Xc};  
  {\it lihat juga} ideal-$\sigma$ ({\bf 112Db})











\indexheader{idempoten}


\indexheader{identik}
variabel acak berdistribusi identik {\bf 273I}, 273Xi,
274I, 274Xj, 276Xd, 276Yg, 285Xq, 285Yc;  
  {\it lihat juga} barisan yang dapat dipertukarkan ({\bf 276Xg})




\indexheader{citra}
filter citra 212Xi, {\bf 2A1Ib}, 2A3Qb, 2A3S


ukuran citra {\bf 112Xf}, 123Ya, 132Yb,
  212Xf, 212Xi,
234C-234F ({\bf 234D}), 
234Xb-234Xd, 
234Ya, 234Yc, 234Yd, 235J, 235Xl,
235Ya, 254Oa, 256G


malapetaka ukuran citra 235H




























\indexheader{tak tentu}


integral tak tentu 131Xa,
  222D-222F, 
222H, {\it 222I}, 222Xa-222Xc, 
222Yc, 224Xg, 225E, {\it 225Od}, 225Xf, {\bf 232D}, 232E, 232Yf, 232Yi;
  {\it lihat juga }
ukuran integral tak tentu
  

ukuran integral tak tentu 234I-234O ({\bf 234J}),
  
234Xh-234Xk, 
234Yi, 234Yj, 234Yl-234Yn, 
235K, {\it 235N}, 235Xh, 235Xl, 235Xm,
253I, 256E, 256J, 256L, 256Xe, 256Ye, 257F, 257Xe, 263Ya, 275Yj, 275Yk,
285Dd, 285Xf, 285Ya;
  {\it lihat juga} ukuran integral tak tentu yang belum dilengkapi


\indexheader{kebebasan}
kebebasan \S272 ({\bf 272A})




\indexheader{independen}






variabel acak bebas {\bf 272Ac},
272D-272I, 
272L, 272M, 272P, 272R-272W, 
272Xb, 272Xd, 272Xe, 272Xh-272Xl, 
272Ya, 272Yb, 272Yd-272Yh, 
273B, 273D, 273E, 273H, 273I,
273L-273N, 
273Xi, 273Xk, 273Xm, 274B-274D, 
274F-274K, 
274Xd, 274Xe, 274Xh, 274Xk, 274Ya, 274Yg,
{\it 275B}, 275Yi, 275Ym, 275Yn, {\it 276Af}, 285I, 285Xi, 285Xj,
285Xp-285Xr, 
285Yc, 285Ym, 285Yn, 285Yt



himpunan bebas {\bf 272Aa}, {\it 272Bb}, 272F, 272N, 273F,
273K, 273Xo




aljabar-$\sigma$ bebas {\bf 272Ab}, 272B, 272D, 272F, 272J,
272K, 272M, 272O, 272Q, 275Ya





\indexheader{indikator}
fungsi indikator (suatu himpunan) {\bf 122Aa}







\indexheader{terinduksi}


topologi terinduksi {\it lihat} topologi subruang ({\bf 2A3C})


\indexheader{induktif}
definisi induktif 2A1B








\indexheader{ketaksamaan}
ketaksamaan {\it lihat} Cauchy, Chebyshev, H\"older, isodiametrik, Jensen, Wirtinger, Young














\indexheader{tak hingga}
tak hingga 112B, 133A, \S135

\indexheader{informasi}




\indexheader{awal}
ordinal awal {\bf 2A1E}, 2A1F, 2A1K




\indexheader{injektif}






\indexheader{dalam}


ukuran dalam 113Yg,
  212Ya, 213Yd
  

----- ----- yang didefinisikan oleh suatu ukuran 113Yh,
  213Xe 


ruang hasil kali dalam 244N, 244Yn, 253Xe




\indexheader{regular dalam}
ukuran regular dalam {\bf 256Ac}





----- ----- terhadap himpunan tertutup 134Xe,
  256Bb, 256Ya


----- ----- terhadap himpunan kompak 134Fb,
  256Bc








\indexheader{terintegralkan}
fungsi terintegralkan \S122 ({\bf 122M}), {\it 123Ya}, {\bf 133B}, {\bf
133Db}, 133F, 133Jd, 133Xa, {\it 135Fa},
  212Bc, 212Fb, 213B, 213G
  {\it lihat juga} fungsi terintegralkan Bochner ({\bf 253Yf}),
$\eusm L^1(\mu)$ ({\bf 242A})


----- {\it lihat juga} terintegralkan secara seragam ({\bf 246A})


\indexheader{integral}
integral (terhadap suatu ukuran)
\S122 ({\bf 122E}, {\bf 122K}, {\bf 122M}), {\bf 133B}, {\bf
133D}; 
  {\it lihat juga}
fungsi terintegralkan, integral Lebesgue ({\bf 122Nb}),
integral bawah ({\bf 133I}), integral Riemann ({\bf 134K}), integral atas ({\bf 133I})










\indexiiheader{integrasi}
integrasi parsial 225F, {\it 225Oe}, 252Xb


integrasi dengan substitusi {\it lihat}
perubahan variabel dalam integrasi




\indexheader{interior}
interior suatu himpunan {\bf 2A3D}, 2A3Ka


\indexiiheader{interpolasi}
interpolasi {\it lihat} Teorema Kekonveksan Riesz







\indexheader{irisan}


\indexheader{interval}
interval
  {\it lihat}
interval setengah terbuka
({\bf 114Aa}, {\bf 115Ab}), interval terbuka ({\bf 111Xb}



\indexheader{invarian}


















 

\indexheader{balik}



transformasi Fourier invers {\bf 283Ab}, 283B, {\bf 283Wa}, 283Xb,
{\bf 284I};
  {\it lihat juga} transformasi Fourier


citra balik (dari suatu himpunan oleh fungsi atau relasi) {\bf 1A1B}

\indexheader{invers pelestari ukuran}

fungsi pelestari ukuran melalui citra balik
134Yl-134Yn, 
  {\bf 234A}, 234B, 234Ea, 234F, 234Xa, 234Xn, 235G, 235Xm, 241Xg,
242Xd, 243Xn, 244Xo, 246Xf, 251L, 251Wp, 251Yb,
254G, 254H, 254Ka, 254O, 254Xc-254Xf, 
254Xh, 254Yc, 254Ye, 256Yd; 
  {\it lihat juga }
ukuran citra ({\bf 234D})




\indexheader{inversi}
Teorema Inversi (untuk deret dan transformasi Fourier)
282G-282J, 
282L, 282O, 282P, 283F, 283I, 283J, 283L, 283Wc, 283Wk, 283Xm,
284C, 284M;
  {\it lihat juga} Teorema Carleson




\indexheader{dapat diinvers}


\indexheader{involusi}












\indexheader{isodiametrik}
ketaksamaan isodiametrik 264H


\indexheader{terisolasi}












\indexiiheader{isomorfisme}
isomorfisme {\it lihat} 
  isomorfisme ruang ukur










\indexmedskip

\indexheader{Jacobian}
Jacobian {\bf 263Ea}







\indexheader{Jensen}




ketaksamaan Jensen 233I, 233J, 233Yi, 233Yj, 242Yi






\indexiiheader{gabungan}
distribusi bersama {\it lihat} distribusi ({\bf 271C})


\indexiiheader{Jordan}


dekomposisi Jordan dari fungsional aditif 231F, 231Ya,
{\it 232C}



\indexmedskip

















\indexheader{Kawada}




















\indexiiheader{kernel}
kernel 
  {\it lihat}
kernel Dirichlet ({\bf 282D}), kernel Fej\'er ({\bf 282D}),
kernel Dirichlet termodifikasi ({\bf 282Xc}),
kernel Poisson ({\bf 282Yg})






\indexiiheader{Kirszbraun}
Teorema Kirszbraun 262C









\indexiiheader{Kolmogorov}


Hukum Kuat Bilangan Besar Kolmogorov 273I, 275Yq







\indexiiheader{Koml\'os}
Teorema Koml\'os 276H, 276Yh


\indexiiheader{K\"onig}
K\"onig, H.\ 232Yk








\indexheader{\Krein}




\indexheader{Kronecker}
Lema Kronecker 273Cb


\indexheader{Kuratowski}















\indexheader{Kwapien}




\indexmedskip

\indexheader{Lacey}
Lema Lacey-Thiele 286M


\indexheader{Laplace}
Teorema Limit Pusat Laplace 274Xf


transformasi Laplace {\bf 123Xc}, {\it 123Yb}, {\bf 133Xd},
  225Xn








\indexheader{secara lateral}


\indexheader{kisi}
kisi {\bf 2A1Ad}; 

  {\it lihat juga}
kisi Banach ({\bf 242G})




----- norma {\it lihat} norma Riesz ({\bf 242Xg})




\indexheader{hukum}
hukum suatu variabel acak {\it lihat} distribusi ({\bf 271C})


hukum bilangan besar {\it lihat} hukum kuat (\S273)


hukum kejadian langka 285Q




\indexheader{Lebesgue}
Lebesgue, H.\   Jilid\ 1 {\it pendahuluan}, bab\ 27 {\it pendahuluan}

Lema Selimut Lebesgue {\it 2A2Ed}


dekomposisi Lebesgue dari fungsional aditif terhitung 232I, 232Yb, 232Yg,
{\it 256Ib}


dekomposisi Lebesgue dari fungsi bervariasi terbatas {\bf
226C}, 226Dc, {\it 226Yb}, 232Yb, 263Ye




Teorema Kerapatan Lebesgue \S223, 261C, {\it 275Xg}


Teorema Konvergensi Terdominasi Lebesgue 123C, 133G

perluasan Lebesgue {\it lihat} pelengkapan ({\bf 212C})


fungsi terintegralkan secara Lebesgue {\bf 122Nb}, 122Yb, 122Ye, {\it
122Yf}

integral Lebesgue {\bf 122Nb}

fungsi terukur Lebesgue {\bf 121C}, 121D, 134Xg, 134Xj,
  {\it 225H}, 233Yd, 262K, 262P, 262Yc
  

himpunan terukur Lebesgue {\bf 114E}, 114F, 114G, {\it 114Xe}, 114Ye, {\bf
115E}, 115F, 115G, 225Gb

ukuran Lebesgue (pada $\Bbb R$) \S114 ({\bf 114E}), 131Xb, 133Xd, 133Xe,
134G-134L, 
  212Xc, {\it 216A}, bab\ 22, 242Xi, 246Yd, 246Ye,
{\it 252N}, {\it 252O},
{\it 252Xf}, {\it 252Xg}, {\it 252Ye}, {\it 252Yo}, \S255


----- ----- (pada $\BbbR^r$) \S115 ({\bf 115E}), 132C, {\it 132Ef}, 133Yc,
\S134,
  211M,  212Xd, 245Yj, 251N, 251Wi, 252Q, 252Xi, 252Yi, 254Xl,
  255A, 255L, 255Xj, 255Ye, 255Yf, 256Ha, 256J, 256L,
264H, 264I, 266C


----- ----- (pada $[0,1]$, $\coint{0,1}$) 211Q, {\it 216Ab},
234Ya, 252Yp, 254K, 254Xh, 254Xk-254Xm, 
254Ye


----- ----- (pada subhimpunan lain dari $\BbbR^r$) {\bf 131B},
  234Ya, 234Yc, 242O, 244Hb, 244I, 244Yi,
246Yf, {\it 246Yl},
251R, 252Yh, 255M-255O, 
255Yg, 255Ym





himpunan terabaikan Lebesgue {\bf 114E}, {\bf 115E}, 134Yk



ukuran luar Lebesgue {\bf 114C}, 114D, 114Xc, {\it 114Yd}, {\bf
115C}, 115D, 115E, 115Xb, 115Xd, 115Yb, 132C, 134A, 134D, 134Fa

himpunan Lebesgue suatu fungsi {\bf 223D},
223Xh-223Xj, 
223Yg, 223Yh, {\bf 261E}, 261Ye, 282Yg


ukuran Lebesgue-Stieltjes {\bf 114Xa}, {\it 114Xb},
114Yb, 114Yc, 114Yf, 131Xc, 132Xg, 134Xd,
  211Xb, 212Xd, 212Xg, 225Xd, 232Xb, {\it 232Yb}, 235Xb, 235Xf, 235Xg,
252Xb, 256Xg, 271Xb, {\it 224Yh}



\indexheader{kiri}








\indexheader{panjang}
panjang kurva 264Yl, 265Xd, 265Ya




panjang suatu interval {\bf 114Ab}




\indexheader{Levi}
Teorema B. Levi {\bf 123A}, 123Xa, 133Ka, 135Ga, 135Hb,
  226E, 242E

sifat Levi pada ruang Riesz bernorma 242Yb, 244Yf


\indexheader{L\'evy}






Teorema konvergensi martingal L\'evy 275I


metrik L\'evy {\bf 274Yc}, {\it 285Yd}


\indexheader{leksikografis}




\indexheader{Liapounoff}
Teorema limit pusat Liapounoff\footnote{A.M.Lyapunov, 1857-1918}
274Xh






\indexheader{pengangkatan}














\indexheader{limit}


limit filter {\bf 2A3Q}, 2A3R, 2A3S


limit barisan {\bf 2A3M}, 2A3Sg




ordinal limit {\bf 2A1Dd}


\indexheader{Lindeberg}
Teorema Limit Pusat Lindeberg 274F-274H, 
285Yo


kondisi Lindeberg {\bf 274H}


\indexheader{Lindel\"of}






\indexiiheader{linear}




\ifdim\pagewidth>467pt\fi
operator linear {\it 262Gc}, 263A, 265B, 265C, {\it \S2A6};  
  {\it lihat juga} operator linear berbatas ({\bf 2A4F}),
operator linear kontinu  

\fontdimen3\tenrm=1.67pt

urutan linear {\it lihat} himpunan terurut total ({\bf 2A1Ac})




topologi ruang linear {\it lihat} ruang topologis linear ({\bf 2A5A}), topologi lemah ({\bf 2A5I}), topologi lemah* ({\bf 2A5Ig})


subruang linear (dari ruang bernorma) 2A4C;
  
  (dari suatu ruang topologis linear) 2A5Ec


ruang topologis linear 245D, {\it 284Ye}, {\bf 2A5A}, 2A5B, 2A5C,
2A5E-2A5I



\indexheader{tertaut}






\indexheader{Lipschitz}
konstanta Lipschitz {\bf 262A}, 262C, 262Ya


fungsi Lipschitz 224Ye, {\it 225Yc}, {\bf 262A},
262B-262E, 
262N, 262Q, 262Xa-262Xc, 
262Xh, 262Xl, 262Ya, 263F, 264G, {\bf 264Yj}; 
  {\it lihat juga} kontinu H\"older ({\bf 282Xj})




\indexiiheader{lokal}
konvergensi lokal dalam ukuran {\it lihat} konvergensi dalam ukuran ({\bf 245A})








\indexiiheader{dapat dilokalkan}




ukuran (ruang ukur) yang dapat dilokalkan {\bf 211G}, 211L, {\it 211Ya}, 211Yb,
212Ga, 213Hb, 213L-213N, 
213Xm, 214Ie, 214K, 214Xa, 214Xc, 214Xe, 215E, 216C, 216E,
{\it 216Ya}, 216Yb, 234Nc, 234O,
234Yk-234Yn, 
241G, 241Ya, 243Gb, 243Hb,
245Ec, 245Yf,
252Yp, {\it 252Yq}, 252Ys, 254U; 
  {\it lihat juga} dapat dilokalkan secara ketat ({\bf 211E})










\indexiiheader{secara lokal}










ukuran (ruang ukur) yang ditentukan secara lokal {\bf 211H}, 211L, {\it 211Ya},
214Xb, {\it 216Xb}, {\it 216Ya}, 216Yb, 251Xd, 252Ya; 
  {\it lihat juga} ukuran lengkap yang ditentukan secara lokal


himpunan terabaikan yang ditentukan secara lokal {\bf 213I},
213J-213L, 
213Xh, 213Xi, 213Xl,
214Ic, 214Xf, 214Xg, 215E, {\it 216Yb}, 234Yl, 252E, 252Yb




ukuran berhingga lokal {\bf 256Ab}, {\it 256C}, 256G, 256Xa, {\it
256Ya}






fungsi terintegralkan lokal 242Xi, {\bf 255Xe}, 255Xf, {\bf 256E},
261E, 261Xa, 262Yh








\indexheader{Lo{\grv e}ve}
Lo{\grv e}ve, M.\  bab\ 27 {\it pendahuluan}


\indexheader{logistik}


\indexheader{Loomis}


\indexheader{Lorentz}




\indexheader{bawah}






integral bawah {\bf 133I}, 133J, 133Xf, 133Yd, 135Ha,
  214Jb, 235A

kerapatan Lebesgue bawah 223Yf




integral Riemann bawah {\bf 134Ka}

fungsi semikontinu bawah {\bf 225H}, 225I, 225Xj, 225Xk,
225Ye, 225Yf







\indexheader{Lusin}




Teorema Lusin 134Yd, 256F
  

\indexheader{Luxemburg}




\indexmedskip




























































\indexiiheader{Markov}




waktu Markov {\it lihat} waktu henti ({\bf 275L}) 













\indexheader{martingal}
martingal \S275 ({\bf 275A}, {\bf 275Cc}, {\bf 275Cd}, {\bf
275Ce});
  {\it lihat juga} martingal balik


teorema konvergensi martingal 275G-275I, 
275K, 275Xf


barisan selisih martingal {\bf 276A}, 276B, 276C, 276E, 276Xe,
276Xf, 276Ya, 276Yb, 276Ye, 276Yg


ketaksamaan martingal 275D, 275F, 275Xb,
275Yd-275Yf, 
276Xb





\indexheader{aliran maksimum}


\indexheader{maksimal}
elemen maksimal dalam himpunan terurut parsial {\bf 2A1Ab}








ketaksamaan maksimal {\it lihat} ketaksamaan maksimal Doob


teorema maksimal 275D, 275Xb, 275Yd, 275Ye, 275Yf,
276Xb, 286A, 286T




\indexheader{Mazur}




\indexheader{McMillan}


\indexheader{McShane}




\indexheader{kurus}




\indexheader{rataan}
rata-rata (variabel acak) {\it lihat} ekspektasi ({\bf 271Ab})




{Teorema Ergodik Rata-Rata }
  {\it lihat} konvergensi dalam rata-rata ({\bf 245Ib})


\indexheader{terukur}






penutup terukur {\it lihat} selubung terukur ({\bf 132D})

selubung terukur {\bf 132D}, 132E, 132F, 132Xg, 134Fc, 134Xd,
  213L, 213M, 214G, 216Yc, 266Xa;  
  {\it lihat juga }
  ukuran luar penuh ({\bf 132F})

sifat selubung terukur {\bf 213Xl}, 214Xk


fungsi terukur (bernilai dalam $\Bbb R$) \S121 ({\bf 121C}),
122Ya,
  212B, 212Fa, 213Ye, 214Ma, 214Na, 235C, 235I, 252O, 252P, 256F, 256Yb,
256Yc

----- -----  (bernilai dalam $\BbbR^r$) {\bf 121Yd},
  256G
  

----- -----  (bernilai dalam ruang lain) {\bf 133Da}, 133E, 133Yb,
{\bf 135E}, {\it 135Xd}, 135Yf  

----- -----  ($(\Sigma,\Tau)$-terukur) {\bf 121Yc},
  235Xc, 251Ya, 251Yd
  

----- -----  {\it lihat juga} terukur Borel, terukur Lebesgue






himpunan terukur {\bf 112A};
  $\mu$-himpunan terukur {\bf 212Cd};
  {\it lihat juga }
terukur secara relatif ({\bf 121A})


ruang terukur {\bf 111Bc}



transformasi terukur \S235; 
  {\it lihat juga} \imp\ fungsi

\indexheader{ukuran}
ukuran {\bf 112A}

----- (dalam `$\mu$ mengukur $E$', `$E$ diukur oleh $\mu$') {\bf $\pmb{>}$112Be}

\indexheader{aljabar ukuran}
aljabar ukuran 211Yb, 211Yc






































----- ----- fungsi {\it lihat} fungsi yang inversnya melestarikan ukuran ({\bf 234A}),
automorfisme ruang ukur, isomorfisme ruang ukur






\indexheader{ruang ukur}
ruang ukur \S112 ({\bf 112A})

automorfisme ruang ukur 254J,
255A, 255L-255N, 
{\it 255 catatan}


isomorfisme ruang ukur 254K,
254Xk-254Xm









\indexiiheader{median}
fungsi median {\bf 2A1Ac}





\indexheader{metrik}
metrik {\bf 2A3F}, 2A4Fb;
  
  metrik Euklides ({\bf 2A3Fb}),
  metrik Hausdorff ({\bf 246Yb}),
  metrik L\'evy ({\bf 274Yc}),
  pseudometrik {\bf 2A3F})




ukuran luar metrik 264Xb, {\bf 264Yc}


ruang metrik {\it 224Ye}, {\it 261Yi}; 
  {\it lihat juga} ruang metrik lengkap, ruang yang dapat dimetrikkan ({\bf 2A3Ff})


\indexheader{secara metrik}






\indexiiheader{dapat dimetrikkan}
ruang topologis yang dapat dimetrikkan {\bf 2A3Ff}, 2A3L; 
  {\it lihat juga} ruang metrik,
ruang yang dapat dimetrikkan dan separabel













\indexiiheader{konveks-tengah}
fungsi konveks-tengah {\bf 233Ya}, 233Yd


\indexiiheader{minimal}
elemen minimal dalam himpunan terurut parsial {\bf 2A1Ab}


\indexiiheader{Minkowski}
ketaksamaan Minkowski {\bf 244F}, 244Yc, 244Ym;
  {\it lihat juga} ketaksamaan Brunn-Minkowski (266C)














\indexheader{termoderasi}




\indexheader{termodifikasi}
kernel Dirichlet termodifikasi {\bf 282Xc}


\indexheader{modular}




\indexheader{modulasi}
modulasi 284Xd, \S284 {\it catatan}, {\bf 286C}


\indexheader{monokompak}


\indexheader{monoton}
Teorema Kelas Monoton 136B-136D, 
136Xc, 136Xf

Teorema Konvergensi Monoton {\it lihat} Teorema B. Levi (123A)



\indexheader{monoton}
fungsi monoton 121D, 222A, 222C, 222Yb, {\it 224D}



\indexiiheader{Monte}
integrasi Monte Carlo 273J, 273Ya




\indexheader{multilinear}
operator multilinear {\bf 253Xc}

















\indexmedskip

\indexheader{Nachbin}


\indexheader{Nadkarni}





















\indexheader{sempit}


\indexheader{alami}









\indexheader{terabaikan}
himpunan terabaikan {\bf 112D}, 131Ca,
  214Cb, 234Hb; 
  {\it lihat juga} 
himpunan terabaikan Lebesgue ({\bf 114E}, {\bf 115E}), ideal nol ({\bf 112Db})




\indexheader{lingkungan}







\indexheader{jaringan}




\indexheader{Neumann}













\indexheader{Nikod\'ym}
Nikod\'ym {\it lihat} Radon-Nikod\'ym






\indexheader{tidak}


\indexheader{tak-menurun}
barisan himpunan tak menurun {\bf 112Ce}

\indexheader{tak-menaik}
barisan himpunan tak meningkat {\bf 112Cf}

\indexheader{tak terukur}
himpunan tak terukur 134B, 134D, 134Xc

\indexheader{tak-utama}













\indexheader{norma}
norma {\bf 2A4B};
  (operator linear) {\bf 2A4F}, 2A4G, 2A4I;
  (matriks) {\bf 262H}, 262Yb;
  (topologi norma) 242Xg, {\bf 2A4Bb};
  {\it lihat juga}
    norma-F ({\bf 2A5B})


\indexheader{normal}
fungsi kerapatan normal 274A, 283N, 283Wi, 283Wj


distribusi normal {\bf 274Ad};
{\it lihat juga} distribusi normal baku ({\bf 274A})


fungsi distribusi normal {\bf 274Aa},
274F-274K, 
274M, 274Xf, 274Xh












variabel acak normal {\bf 274A}, 274B, 274Xb,
285E, 285Xp, 285Xq, 285Xh














\indexheader{ternormalisasi}




ukuran Hausdorff ternormalisasi 264 {\it catatan},
\S265 ({\bf 265A})


\indexheader{penormal}


\indexheader{bernorma}
aljabar bernorma {\bf 2A4J};  
  {\it lihat juga} aljabar Banach ({\bf 2A4Jb})


ruang bernorma {\it 224Yf}, \S2A4 ({\bf 2A4Ba}); 
  {\it lihat juga} ruang Banach ({\bf 2A4D})












\indexiiheader{tak di mana pun}
ukuran yang pada bagian mana pun tidak mengukur semua himpunan {\bf 214Yd}
















\indexheader{nol}
ideal nol ukuran {\bf 112Db},
  251Xp




himpunan nol {\it lihat} terabaikan ({\bf 112Da})


\indexmedskip

\indexheader{ganjil}
fungsi ganjil 255Xb, 283Ye







\indexheader{Ogasawara}
















\indexheader{terbuka}
selang terbuka {\it 111Xb}, 114G, 115G, 1A1A,
  2A2I
  



himpunan terbuka (dalam $\BbbR^r$) {\bf 111Gc}, {\it 111Yc}, {\it 114Yd}, {\it 115G},
{\it 115Yb}, {\bf 133Xc}, 134Fa, 134Xe,
{\bf 135Xa}, {\bf 1A2A}, 1A2B, 1A2D;
  (dalam $\Bbb R$) 111Gc, 111Ye, {\it 114G}, 134Xd, 2A2I; (dalam ruang topologis lain) 256Yf, {\bf 2A3A}, 2A3G;
  {\it lihat juga} topologi ({\bf 2A3A})








\indexiiheader{opsional}
waktu opsional {\it lihat} waktu henti ({\bf 275L}) 












\indexiiheader{berbatas orde}


himpunan berbatas orde (dalam ruang terurut parsial) {\bf 2A1Ab}















\indexiiheader{lengkap orde}
lengkap orde {\it lihat} lengkap Dedekind ({\bf 241Ec})


\indexiiheader{kontinu orde}






norma kontinu orde (pada ruang Riesz) {\bf 242Yc}, 242Ye,
244Ye












\indexiiheader{konvergen menurut urutan*}
barisan konvergen menurut urutan*; (dalam $L^0(\mu)$) {\bf 245C}, 245K, 245L, 245Xc, 245Xd
































\indexiiheader{satuan orde}
satuan orde (dalam ruang Riesz) 243C




\indexheader{terurut}


himpunan terurut {\it lihat} himpunan terurut parsial ({\bf 2A1Aa}), himpunan terurut total ({\bf 2A1Ac}), himpunan terurut baik ({\bf 2A1Ae})


\indexiiheader{pengurutan}
pengurutan ukuran {\it 214Hb}, 
{\bf 234P}, 234Q, 234Xl, 234Yo


\indexiiheader{ordinal}
ordinal {\bf 2A1C}, 2A1D-2A1F, 
2A1K










\indexiiheader{ordinat}
himpunan ordinat {\bf 252N}, 252Yl, 252Ym, 252Yv










\indexiiheader{ortogonal}




matriks ortogonal {\bf 2A6B}, 2A6C


proyeksi ortogonal dalam ruang Hilbert 244Nb, 244Yk, 244Yl




\indexiiheader{ortonormal}


vektor ortonormal {\bf 2A6B}












\indexheader{luar}




\indexheader{ukuran luar}
ukuran luar \S113 ({\bf 113A}), 114Xd, {\bf 132B}, {\it 132Xg},
136Ya,
  {\it 212Xf}, 213Xa, 213Xg, 213Xi, 213Yb, 251B, 251Wa, {\it 251Xe},
254B, {\it 264B}, {\it 264Xa}, {\it 264Ya}, {\it 264Yo};  
  {\it lihat juga}
ukuran luar Lebesgue ({\bf 114C}, {\bf 115C}), ukuran luar metrik ({\bf 264Yc}),
  ukuran luar regular ({\bf 132Xa})

----- ----- yang didefinisikan dari suatu ukuran 113Ya,
132A-132E ({\bf 132B}), 
132Xa-132Xi, 
132Xk, 132Ya-132Yc, 
132Yg, 133Je,
  212Ea, {\it 212Xa}, 212Xb, 213C, 213Fb, 213Xa, 213Xg, 213Xi, 213Xj,
213Yb, 213Ye, 214Cd, {\it 215Yc}, 234Bf, 234Ya,
251P, 251S, 251Wk, 251Wm, 251Xn,
251Xq, {\it 252D}, {\it 252I}, {\it 252O}, 252Ym,
{\it 254G}, 254L, 254S, 254Xb, 254Xr, 254Yf, 256Xi, 264Fb, 264Ye



\indexiiheader{regular luar}


ukuran regular luar 134Fa, 134Xe, 256Xi









\indexmedskip













\indexiiheader{Parseval}
identitas Parseval {\bf 284Qd};  {\it lihat juga} Teorema Plancherel


\indexheader{parsial}
turunan parsial {\it 123D},
  {\it 252Ye}, 262I, 262J, 262Xh, 262Yb, 262Yc 



urutan parsial
  {\it lihat} himpunan terurut parsial ({\bf 2A1Aa})




\indexheader{terurut parsial}
ruang linear terurut parsial {\bf 241E}, 241Yg


himpunan terurut parsial 234Qa, 234Yo, {\bf 2A1Aa}


\indexheader{partisi}
partisi {\bf 1A1J}





\indexheader{Peano}
kurva Peano 134Yl-134Yo;   
  {\it lihat juga} kurva pengisi ruang

\indexheader{sempurna}








\indexheader{secara sempurna}


\indexheader{keliling}




\indexheader{periodik}


perluasan periodik fungsi pada $\ocint{-\pi,\pi}$ {\bf 282Ae}












\indexheader{Perron}


\indexheader{Pettis}






\indexheader{Pfeffer}


















\indexheader{Plancherel}
Teorema Plancherel (untuk deret dan transformasi Fourier fungsi terintegralkan-kuadrat) 282K, 284O, 284Qd







\indexheader{Poincar\'e}


\indexheader{terhitung-titik}


\indexheader{berhingga-titik}




\indexheader{bertumpu pada titik}
ukuran yang bertumpu pada titik {\bf 112Bd},
  {\bf 211K}, 211O, 211Qb, {\it 211Rc}, 211Xb, 211Xf, 213Xo,
234Xb, 234Xe, 234Xk, 251Xu, 256Hb;
  {\it lihat juga} ukuran Dirac ({\bf 112Bd}), ukuran empiris


\indexiiheader{titik demi titik}


konvergensi titik demi titik (topologi pada ruang fungsi)
281Yf






konvergen titik demi titik 
  {\it lihat} konvergen menurut urutan* ({\bf 245Cb})


topologi titik demi titik {\it lihat} konvergensi titik demi titik


\indexheader{Poisson}
distribusi Poisson 285Q, {\bf 285Xr}


kernel Poisson {\bf 282Yg}




Teorema Poisson 285Q




\indexheader{polar}
koordinat polar 263G, 263Xf









\indexheader{P\'olya}
skema urna P\'olya 275Xc


\indexheader{polinomial}
polinomial (pada $\BbbR^r$) 252Yi








\indexheader{berpori}
himpunan berpori {\bf 223Ye}, {\bf 261Yg}, 262L


\indexheader{positif}
kerucut positif 253Gb, 253Xi, 253Yd


fungsi definit positif 283Xs, 285Xu













\indexiiheader{himpunan kuasa}
himpunan kuasa $\Cal PX$
  (ukuran biasa pada) {\bf 254J}, 254Xf, 254Xr, 254Yf




{----- $\Cal P\Bbb N$ }
1A1Hb,
  2A1Ha, 2A1Lb;  
  (ukuran biasa pada)  273G, 273Xe, 273Xf
  



















\indexheader{dapat diprediksi}
barisan terprediksi {\bf 276Ec}





\indexheader{pra-Radon}


\indexheader{hampir}
hampir di mana-mana {\bf 112De}




\indexheader{primitif}
ukuran produk primitif {\bf 251C}, 251E, 251F, 251H, 251K, 251Wa,
251Xb-251Xh, 
251Xk, 251Xm-251Xo, 
251Xq, 251Xr, 252Yc, 252Yd, 252Yl, 253Ya-253Yc, 
253Yg


\indexheader{utama}






ultrafilter utama {\bf 2A1N}


\indexheader{peluang}








fungsi kerapatan peluang {\it lihat} fungsi kerapatan ({\bf 271H}),
fungsi kerapatan normal ({\bf 274A})


ruang peluang {\bf 211B}, 211L, 211Q, {\it 211Xb}, {\it 211Xc}, {\it 211Xd},
{\it 212Ga}, {\it 213Ha}, 215B, 234Bb, {\it 243Xi},
{\it 245Xe}, 253Xh, \S254, bab\ 27






\indexheader{produk}





  





ukuran produk bab\ 25; 
  {\it lihat juga} versi c.l.d.\ ukuran produk ({\bf 251F}, {\bf 251W}),
ukuran produk primitif ({\bf 251C}),
ukuran peluang produk ({\bf 254C})












ukuran peluang produk \S254 ({\bf 254C}), 272G, 272J, 272M,
273J, 273Xj, 275J, 275Yj, 275Yk, {\it 281Yk}; 
  {\it lihat juga } $\{0,1\}^I$















topologi produk 281Yc, {\bf 2A3T}




----- {\it lihat juga} ruang hasil kali dalam






















































\indexiiheader{pseudometrik}
pseudometrik {\bf 2A3F},
2A3G-2A3M 
2A3S-2A3U, 
{\it 2A5B}





\indexheader{pseudo-sederhana}
fungsi pseudo-sederhana {\bf 122Ye}, 133Ye



\indexheader{tarikan-balik}
ukuran tarik-balik 234F, 234Ye


\indexheader{secara murni}


ukuran atomik murni (ruang) {\bf 211K}, 211N, 211R, {\it 211Xb, 211Xc, 211Xd},
{\it 212Gd}, 213He, 214K, 214Xd, 234Be, 234Xe, 234Xj, 234Xn, 251Xu


himpunan terukur tak hingga murni {\bf 213 catatan}




\indexiiheader{dorong}
ukuran dorong-maju
{\it lihat} ukuran citra ({\bf 234D})




\indexmedskip

\indexheader{kuasi-diadik}


\indexheader{kuasi-homogen}











\indexheader{kuasi-rapat-orde}


\indexheader{kuasi-Radon}
ukuran kuasi-Radon (ruang) 256Ya, {\it 263Ya}




\indexheader{kuasi-sederhana}
fungsi kuasi-sederhana {\bf 122Yd}, 133Yd










\indexheader{hasil bagi}






ruang linear terurut parsial hasil bagi 241Yg;
  {\it lihat juga} ruang Riesz hasil bagi


ruang Riesz hasil bagi 241Yg, 241Yh, 242Yg






topologi hasil bagi {\it 245Ba}


\indexmedskip

\indexiiheader{Rademacher}
Teorema Rademacher 262Q


\indexiiheader{Radon}


\indexiiheader{ukuran Radon}


{ukuran Radon}
(pada $\Bbb R$ atau $\BbbR^r$)
\S256 ({\bf 256A}), 257A, {\it 284R}, 284Yi 


----- (pada $\ocint{-\pi,\pi}$) {\bf 257Yb}




turunan Radon--Nikod\'ym {\bf 232Hf}, 232Yj,
234J, 234Ka, 234Yi, 234Yj, 235Xh, 256J, 257F, 257Xe,
272Xe, 275Yb, 275Yj,
285Dd, 285Xf, 285Ya


Teorema Radon--Nikod\'ym 232E-232G, 
234O, 235Xj, 242I, {\it 244Yl}


ukuran peluang Radon (pada $\Bbb R$ atau $\BbbR^r$) 256Hc, 257A,
257Xa, 271B, 271C, 271Xb, 274Xi, 274Yd,
285A, 285D, 285F, 285J, 285M, 285R, 285Xa, 285Xf, 285Xg, 285Xk, 285Xm,
285Xn, 285Xp, 285Ya, 285Yb, 285Yh, 285Yp; 


----- ----- ----- (pada ruang lain) {\bf 256Yf}, 271Ya


ukuran produk Radon (dari sejumlah hingga ruang) 256K











\indexheader{acak}




variabel acak {\bf 271Aa}














\indexheader{menurun cepat}
fungsi uji menurun cepat \S284 ({\bf 284A}, {\bf 284Wa}),
285Dc, 285Xe, 285Ya




















\indexheader{penyusunan ulang}
penyusunan ulang {\it lihat} penyusunan ulang menurun ({\bf 252Yo})














\indexheader{rekuren}


\indexheader{rekursi}
rekursi 2A1B





































\indexheader{regular}








ukuran regular {\it lihat} 
regular dalam ({\bf 256Ac})


















ukuran luar reguler 132C, {\bf 132Xa},
  {\it 213C}, {\it 213Xa}, 214Hb, 214Xb, 251Xn, 254Xb, 264Fb
  



ruang topologis regular {\bf 2A5J}


\indexheader{secara reguler}






\indexheader{relasi}
relasi 1A1B












\indexheader{secara relatif}


himpunan relatif kompak $\pmb{>}${\bf 2A3Na}, 2A3Ob; 
  {\it lihat juga} relatif kompak lemah




















himpunan terukur secara relatif {\bf 121A}



himpunan relatif kompak lemah (dalam suatu ruang bernorma) 247C,
{\bf 2A5Id}









\indexheader{berulang}
integral berulang \S252 ({\bf 252A}); 
  {\it lihat juga} Teorema Fubini, Teorema Tonelli




\indexheader{representasi}



















\indexheader{mempertahankan}




\indexheader{retrak}




\indexheader{balik}


martingal terbalik 275K





\indexheader{Riemann}


kriteria Riemann {\it 281Yh}


fungsi terintegralkan Riemann {\bf 134K}, 134L,
  281Yh, 281Yi

integral Riemann {\bf 134K},
  242 {\it catatan}
  

Lema Riemann-Lebesgue 282E, 282F




\indexheader{Riesz}
Teorema Kekonveksan Riesz 244 {\it catatan}






norma Riesz {\bf 242Xg}






ruang Riesz (= kekisi vektor) 231Yc, {\bf 241Ed}, 241F,
241Yc, 241Yg; 
  {\it lihat juga}
kisi Banach ({\bf 242G}),
norma Riesz ({\bf 242Xg})





















































\indexmedskip

\indexheader{Saks}
Saks {\it lihat} Teorema Denjoy-Young-Saks (222L),
Teorema Vitali-Hahn-Saks (246Yi)






\indexheader{saltus}
fungsi saltus {\bf 226B}, {\bf 226Db}, 226Xb


bagian saltus fungsi bervariasi terbatas {\bf 226C},
226Xc, 226Xd, 226Yd




















\indexiiheader{skalar}
perkalian skalar ukuran {\bf 234Xf}, 234Xl


\indexheader{secara skalar}


\indexheader{terpencar}








\indexheader{Schr\"oder}
Teorema Schr\"oder-Bernstein 2A1G


\indexheader{Schwartz}
fungsi Schwartz {\it lihat} fungsi uji menurun cepat
({\bf 284A})


\indexheader{ala Schwartz}
distribusi Schwartz 284R, 284 {\it catatan};
  {\it lihat juga} distribusi tempered ({\bf 284 {\it catatan}})




\indexheader{kedua}














\indexheader{adjoin-diri}


\indexheader{menopang-diri}
himpunan menyangga-diri (dalam suatu ruang ukur topologis)
{\bf 256Xf}







\indexheader{semikompak}


\indexheader{semikontinu}
fungsi semikontinu
  {\it lihat} semikontinu bawah ({\bf 225H})


\indexheader{semi-definit}


\indexheader{semifinit}


ruang ukur semifinit {\bf 211F}, 211L, 211Xf, {\it 211Ya},
{\it 212Ga}, 213A, 213B, 213Hc, 213Xc, 213Xd, 213Xh, 213Xl, 213Xm,
213Ya, 213Yb-213Yd 
214Xg, {\it 214Ya}, {\it 215B},
{\it 216Xa}, {\it 216Yb}, 216Ye, 234B, 234Na, 234Xe, 234Xi,
{\it 235M}, {\it 235Xd}, {\it 241G}, {\it 241Ya}, {\it 241Yd},
243Ga, 245Ea, 245J, {\it 245Xd}, 245Xk, 245Xm, 246Jd,
246Xh, 251J, 251Xd, 252P, {\it 252Yk}, 253Xf, 253Xg


versi semifinit suatu ukuran {\bf 213Xc}, 213Xd





\indexiiheader{seminorma}
seminorma {\bf 2A5D}, 2A5Ia;  
{\it lihat juga} seminorma-F ({\bf 2A5B})





\indexheader{semigelanggang}
semigelanggang himpunan {\bf 115Ye}



\indexiiheader{separabel}
ruang topologis separabel {\bf 2A3Ud}


ruang Banach separabel 244I, 254Yd




ruang dapat dimetrikkan separabel 245Yj, 264Yb, 284Ye





\indexheader{memisahkan}


\indexheader{terpisah}





\indexheader{pemisahan}














\indexheader{secara sekuensial}


































































\indexheader{Sierpi\'nski}
Sierpi\'nski {\it lihat} Teorema Kelas Monoton (136B)



\indexheader{bertanda}
ukuran bertanda 
  {\it lihat} fungsional aditif terhitung ({\bf 231C})











\indexheader{sederhana}
fungsi sederhana \S122 ({\bf $\pmb{>}$122A}), 242M
  











\indexheader{Sina\v\i}


\indexheader{singular}
fungsional aditif singular {\bf 232Ac}, 232I, {\bf 232Yg}




ukuran singular 231Yf, {\bf 232Yg}



\indexheader{miring}




\indexheader{slalom}




\indexheader{kecil}


\indexheader{mulus}


fungsi mulus (pada $\Bbb R$ atau $\BbbR^r$) {\bf 242Xi}, 255Xf,
262Ye-262Yh, 
274Yb, {\bf 284A}, {\bf 284Wa}, 285Xg




\indexheader{penghalusan}
penghalusan dengan konvolusi 261Ye







\indexheader{solid}
selubung solid (dari suatu subhimpunan dari suatu ruang Riesz) 247Xa

































\indexheader{ruang}
kurva pengisi-ruang 134Yl, 254Ye 




\indexheader{spektral}


\indexheader{spektrum}


\indexiiheader{bola}
ukuran permukaan bola pada 265F-265H, 
265Xa-265Xc, 
265Xe, 265Xf




\indexheader{sferis}
koordinat polar sferis 263Xf, 265F














\indexheader{persegi}


fungsi terintegralkan-kuadrat {\bf 244Na};
  {\it lihat juga} $\eusm L^2$




\indexheader{penstabil}







  





distribusi normal baku, variabel acak normal baku {\bf
274A}
















\indexiiheader{Steiner}
simetrisasi Steiner 264H





\indexiiheader{tangga}
fungsi tangga {\bf 226Xb}


\indexheader{Stieltjes}
ukuran Stieltjes {\it lihat}
ukuran Lebesgue-Stieltjes ({\bf 114Xa})


\indexiiheader{Stirling}
rumus Stirling {\bf 252Yu}







\indexheader{secara stokastik}
bebas stokastik {\it lihat} bebas ({\bf 272A})


\indexheader{Stone}












teorema Stone--Weierstrass 281A, 281E, 281G, 281Ya, 281Yg


\indexheader{penghentian}
waktu henti {\bf 275L}, 275M-275O, 
275Xi-275Xk


\indexheader{langsung}


\indexheader{Strassen}


\indexheader{strategi}


\indexiiheader{secara ketat}
ruang ukur yang dapat dilokalkan secara ketat {\bf 211E},
211L, {\it 211N}, 211Xf, {\it 211Ye}, 212Gb,
213Ha, 213J, 213O, 213Xa, 213Xj, 213Xn, 213Yf,
214Ia, 214K, 215Xf, {\it 216E}, 216Yd,
234Nd, {\it 235N}, 251O, 251Q, 251Wl, 251Xo, 252B, 252D,
252Yr, {\it 252Ys}









\indexheader{kuat}
hukum kuat bilangan besar
273D-273J, 
275Yq, 276C, 276F, 276Ye, {\it 276Yg}








































\indexheader{subaljabar}
subaljabar;  
  {\it lihat} 
$\sigma$-subaljabar ({\bf 233A})





\indexheader{subdivisi}


\indexheader{subgrup}





\indexheader{subhomomorfisme}


\indexheader{subkisi}


\indexheader{submartingal}
submartingal {\bf 275Yg}, 275Yh













\indexheader{subruang}
ukuran subruang 113Yb,
  214A, {\bf 214B}, 214C, 214H, 214I, 214Q,
214Xb-214Xg, 
{\it 214Ya}, 216Xa, 216Xb, 241Yg, 242Yf, 243Ya, 244Yd,
245Yb, 251Q, 251R, 251Wl, 251Xo, 251Xp, 251Yc, 254La, 254Yg, 264Yf  

----- -----  pada subhimpunan terukur 131A, {\bf 131B}, 131C, 132Xb,
135I,
  214K, 214L, 214Xa, 214Xh, 241Yf, 247A  

----- -----  (integrasi terhadap ukuran subruang) {\bf 131D},
131E-131H, 
131Xa-131Xc, 
{\it 133Dc}, 133Xb, 135I,
  {\bf 214D}, 214E-214G, 
214N, 214J, 214Xl

subruang dari suatu ruang bernorma 2A4C


topologi subruang {\bf 2A3C}, 2A3J




aljabar-$\sigma$ subruang {\bf 121A},
  {\it 214Ce}

\indexiiheader{substitusi}
substitusi {\it lihat} perubahan variabel dalam integrasi





\indexheader{penerus}
kardinal penerus 2A1Fc


----- ordinal {\bf 2A1Dd}


\indexheader{jumlah}
jumlah atas himpunan indeks sembarang 112Bd, 226A

jumlah ukuran 112Yf, 122Xi,
  {\bf 234G}, 234H, 234Qb, 234Xd-234Xg, 
234Xi, 234Xl, 234Yf-234Yh




\indexheader{dapat dijumlahkan}
keluarga bilangan real yang dapat dijumlahkan {\bf 226A}, 226Xa








\indexheader{supermodular}


\indexheader{dukungan}




----- dari suatu ukuran topologis {\bf 256Xf}, 257Xd, 285Xg




----- {\it lihat juga} dukungan terbatas,
  dukungan kompak


\indexheader{tertumpu}


{tertumpu {\it lihat}}\
tertumpu-titik ({\bf 112Bd})

\indexiiheader{penyangga}
penyangga
  {\it lihat} himpunan menyangga-diri ({\bf 256Xf}), dukungan


\indexheader{supremum}
supremum {\bf 2A1Ab}




\indexiiheader{permukaan}
ukuran permukaan {\it lihat} ukuran Hausdorff ternormalisasi ({\bf 265A})



\indexheader{simetris}
distribusi simetris 272Yc










\indexheader{simetrisasi}
simetrisasi {\it lihat} simetrisasi Steiner






\indexmedskip























\indexheader{bertemper}
distribusi tempered {\bf 284 catatan}


----- fungsi \S284 ({\bf 284D}), {\it 286D}


----- ukuran {\bf 284Yi}


\indexheader{tensor}
hasil kali tensor ruang linear 253 {\it catatan}


\indexheader{tenda}


\indexheader{uji}
fungsi uji 242Xi, 284 {\it catatan\/};
  {\it lihat juga} fungsi uji menurun cepat ({\bf 284A})




\indexheader{tebal}
himpunan tebal {\it lihat} berukuran luar penuh ({\bf 132F})




\indexiiheader{tiga}
Teorema Tiga Deret 275Yn




\indexheader{Tietze}


\indexheader{ketat}
{ketat {\it lihat} ketat seragam ({\bf 285Xm})}

























\indexiiheader{Tonelli}
Teorema Tonelli 252G, 252H, 252R


\indexiiheader{topologis}


ruang ukur topologis {\bf 256Aa}





ruang topologis \S2A3 ({\bf 2A3A})


ruang vektor topologis {\it lihat} ruang topologis linear
({\bf 2A5A})


\indexheader{topologi}
topologi \S2A2, \S2A3 ({\bf 2A3A}); 
  {\it lihat juga} konvergensi dalam ukuran ({\bf 245A}),
topologi ruang linear ({\bf 2A5A})





\indexheader{total}
urutan total
  {\it lihat} himpunan terurut total ({\bf 2A1Ac})


variasi total (dari suatu fungsional aditif) {\bf 231Yh};
(dari suatu fungsi) {\it lihat} variasi ({\bf 224A})






\indexheader{secara total}




ukuran hingga total (ruang) {\bf 211C}, 211L, {\it 211Xb, 211Xc, 211Xd},
{\it 212Ga}, {\it 213Ha}, 214Ia, 214Ka, {\it 214Yc}, {\it 215Yc},
{\it 232Bd}, {\it 232G}, {\it 243Ic}, 234Bc, {\it 243Xk}, {\it 245Fd},
{\it 245Ye}, 246Xi, 246Ya







himpunan terurut total {\it 135Ba}, {\bf
  2A1Ac}
  






\indexheader{jejak}




{jejak }
(dari suatu aljabar-$\sigma$) {\it lihat} aljabar-$\sigma$ subruang ({\bf 121A})



\indexheader{transhingga}
rekursi transfinit 2A1B





\indexheader{transitif}


\indexheader{invarian terhadap translasi}










{invarian terhadap translasi }
ukuran 114Xf, 115Xd, 134A, 134Ye, 134Yf,
  255A, 255N, 255Yn
  





\indexheader{transportasi}


\indexheader{transversal}








\indexheader{secara sejati}
fungsional aditif kontinu sejati {\bf 232Ab},
232B-232E, 
232H, 232I, 232Xa, 232Xb, {\it 232Xf}, 232Xh, 232Ya, 232Ye



\indexheader{terpangkas}






















\indexheader{Tychonoff}


\indexheader{tipe}



\indexmedskip

\indexiiheader{Ulam}
Ulam S.\ {\it lihat} masalah Banach-Ulam







\indexheader{ultrafilter}
ultrafilter  254Yf, {\bf 2A1N}, 2A1O, 2A3R, 2A3Se; 
  {\it lihat juga} ultrafilter utama ({\bf 2A1N})




Teorema Ultrafilter 2A1O




\indexheader{belum dilengkapi}
ukuran integral tak tentu belum dilengkapi {\bf 234Kc}






\indexheader{uniferent}


\indexheader{seragam}












ruang seragam {\it 2A5F}













\indexiiheader{seragam}


fungsi kontinu seragam {\it 224Xa}, 255K, 262A,
{\bf 2A2Cc}




ruang bernorma konveks seragam 244O, 244Yn, {\bf 2A4K}


barisan terdistribusi seragam {\it lihat} terekuidistribusi ({\bf 281Yi}) 











himpunan terintegralkan seragam (dalam $\eusm L^1$)
\S246 ($\pmb{>}${\bf 246A}), 252Yo, 272Ye, 273Na, 274J,
275H, 275Xj, {\it 275Yp}, 276Xe, 276Yb; 
  (dalam $L^1(\mu)$) \S246 ({\bf 246A}), 247C, 247D, 247Xe, 253Xd








ketat seragam 
(himpunan ukuran) {\bf 285Xm}, 285Xn, {\bf 285Yf}, 285Yg









\indexiiheader{satuan}
bola satuan dalam $\BbbR^r$ 252Q










\indexheader{universal}


Teorema Pemetaan Universal 253F, 254G
































\indexheader{penyilangan-naik}
penyilangan-naik 275E, 275F








\indexheader{atas}








integral atas {\bf 133I}, 133J-133L, 
133Xf, 133Yd, {\bf 135H}, 135I, 135Yc,
  212Xj, 214Ja, 214Xl, 235A, 252Yj, 252Ym, 252Yn, 253J, 253K, 273Xj

integral Riemann atas {\bf 134Ka}





\indexheader{ke atas}








himpunan terurut parsial terarah ke atas
{\bf 2A1Ab}
















\indexheader{Urysohn}






\indexheader{biasa}
ukuran biasa pada $\{0,1\}^I$ {\bf 254J};
  {\it lihat di bawah} $\{0,1\}^I$




----- ----- pada $\Cal PX$ {\bf 254J};  {\it lihat di bawah} himpunan kuasa





\indexmedskip

\indexheader{samar}
topologi samar (pada suatu ruang dari ukuran bertanda) {\bf 274Ld}, 274Xi,
274Yc-274Yf, 
275Yl, 285K, 285L, 285S, 285U, 285Xn, 285Xt, 285Yd, 285Ye, 285Yh, 285Yi,
285Yk, 285Yp, 285Yt





\indexiiheader{varians}
varians (suatu distribusi) {\bf 271F}


----- (dari suatu variabel acak) {\bf 271Ac}, 271Xa, 272S,
274Ya, 274Yg, 285Gb, 285Xr


\indexheader{variasi}
variasi dari suatu fungsi \S224 ({\bf 224A}, {\bf 224K}, {\bf 224Yd},
{\bf 224Ye}), 225Yd, 226B, 226Db, 226Xd, 226Xe, 226Yc, 226Yd,
264Xf, 265Yb;
  {\it lihat juga} terbatas variasi ({\bf 224A})


----- dari suatu ukuran {\it lihat} variasi total ({\bf 231Yh})




\indexheader{vektor}
integrasi vektor {\it lihat} Bochner integral ({\bf 253Yf})


kekisi vektor {\it lihat} ruang Riesz ({\bf 241E})






\indexheader{sangat}







\indexheader{secara virtual}
fungsi terukur secara virtual {\bf 122Q}, 122Xe, 122Xf, 135Ia,
  212Bb, 212Fa, 234K, 241A, 252E, 252O

\indexheader{Vitali}
konstruksi Vitali atas suatu himpunan tak terukur 134B

penutup Vitali 261Ya


Teorema Vitali 221A, 221Ya,
221Yc-221Ye, 
261B, {\it 261Yg}, 261Yk


Teorema Vitali-Hahn-Saks 246Yi




\indexheader{volume}
volume {\bf 115Ac}

----- dari suatu bola dalam $\BbbR^r$ 252Q, 252Xi



\indexmedskip

\indexiiheader{Wald}
persamaan Wald 272Xh










\indexiiheader{lemah}










{topologi lemah }





----- ----- {\it lihat juga}
 
(secara relatif) kompak lemah, 
konvergen lemah 




topologi lemah* pada suatu ruang dual 253Yd, 285Yh, {\bf 2A5Ig}; 
  {\it lihat juga} topologi samar ({\bf 274Ld})





\indexheader{secara lemah}




himpunan kompak lemah (dalam suatu ruang topologis linear) 247C, 247Xa,
247Xc, 247Xd, {\bf 2A5I}; 
  {\it lihat juga}
relatif kompak lemah ({\bf 2A5Id})




barisan konvergen lemah dalam suatu ruang bernorma 247Yb



































\indexheader{Weierstrass}
Teorema Aproksimasi Weierstrass 281F;
  {\it lihat juga} Teorema Stone-Weierstrass





\indexiiheader{terdistribusi baik}
barisan terdistribusi baik 281Xh, {\bf 281Ym}








\indexiiheader{terurut baik}
himpunan terurut baik 214P, {\bf 2A1Ae}, 2A1B, 2A1Dg, 2A1Ka; 
  {\it lihat juga} ordinal ({\bf 2A1C})


\indexiiheader{pengurutan baik}
Teorema Pengurutan Baik 2A1Ka








\indexiiheader{Weyl} Teorema Ekuidistribusi Weyl 281M, 281N, 281Xh












\indexheader{Wirtinger}
ketaksamaan Wirtinger 282Yf












\indexmedskip

\indexheader{Young}


ketaksamaan Young 255Yl


----- {\it lihat juga} Teorema Denjoy-Young-Saks (222L)



\indexmedskip

\indexheader{Zermelo}
Teorema Pengurutan Baik Zermelo 2A1Ka




\indexheader{nol}


hukum nol-satu 254S, 272O, 272Xf, 272Xg








\indexheader{Zorn}
Lema Zorn 2A1M



\indexmedskip 

\indexheader{$\scriptstyle\pmb{A}$}















a.e.\ (`hampir di mana-mana') {\bf 112Dd}

a.s.\ (`hampir pasti') {\bf 112De}



\indexmedskip
\indexiiheader{$\scriptstyle\pmb{B}$}



$B$ (dalam $B(x,\delta)$, bola tertutup) {\bf 261A}


$\eurm B$ (dalam $\eurm B(U;V)$, ruang dari operator linear terbatas)
253Xb, 253Yj, 253Yk, {\bf 2A4F}, 2A4G, 2A4H 

















\indexmedskip
\indexheader{$\scriptstyle\pmb{C}$}







$\frak c$ (kardinal dari $\Bbb R$ dan $\Cal P\Bbb N$) 2A1H,
{\bf 2A1L}, 2A1P




$\Bbb C$ = himpunan bilangan kompleks; (dalam $\RoverC$) 2A4A

$C$ (dalam $C(X)$, dengan $X$ suatu ruang topologis) 243Xo, 281Yc,
281Ye, 281Yf;  
  (dalam $C([0,1])$) 242 {\it catatan}




$C_b$ (dalam $C_b(X)$, dengan $X$ suatu ruang topologis) {\bf 281A},
281E, 281G, 281Ya, 281Yd, 281Yg, 285Yh;  
  (dalam $C_b(X;\Bbb C)$) 281G, 281Yg




$C_k$ (dalam $C_k(X)$, dengan $X$ suatu ruang topologis)
242O, 244Hb, 244Yj, 256Xh;  
  {\it lihat juga} dukungan kompak


----- (dalam $C_k(X;\Bbb C)$) 242Pd

































c.l.d.\ {\bf 251F}, 251G,
251I-251L, 
251N-251U, 
{\bf 251W}, 251Xb-251Xm, 
  
251Xp, 251Xs-251Xu, 
251Yb-251Yd, 
\S\S252-253, 254Db, 254U, 254Yg, 256K, 256L


versi c.l.d.\ dari suatu ukuran (ruang) {\bf 213E},
213F-213H, 
213M, 213Xb-213Xe, 
213Xg, 213Xi, 213Xn, 213Xo, 213Yc, {\it 214Xe}, 214Xi, {\it 232Ye},
234Xl, 234Yj, 234Yo, 241Ya, 242Yh, {\it 244Ya},
245Yc, 251Ic, 251T, 251Wf, 251Wl,
251Xe, 251Xk, 251Xl, 252Ya










\indexmedskip
\indexheader{$\scriptstyle\pmb{D}$}









$\DiniD$ (dalam $\DiniD^+f$, $\DiniD^-f$) {\it lihat} turunan Dini
({\bf 222J})


$\Dinid$ (dalam $\Dinid^+f$, $\Dinid^-f$) {\it lihat} turunan Dini
({\bf 222J})














$\diam$ (dalam $\diam A$) = diameter



$\dom$ (dalam $\dom f$):  domain suatu fungsi $f$

\indexmedskip
\indexheader{$\scriptstyle\pmb{E}$}

$\Expn$ (dalam $\Expn(X)$, ekspektasi dari suatu variabel acak)
{\bf 271Ab}



$\esssup$ {\it lihat} supremum esensial ({\bf 243Da})




\indexmedskip
\indexheader{$\scriptstyle\pmb{F}$}





















$f$ subruang 241H, 241 {\it catatan\/}


norma-F {\bf 2A5B}


seminorma-F 245D, {\bf 2A5B}, 2A5C, 2A5D, 2A5G 


\indexmedskip
\indexheader{$\scriptstyle\pmb{G}$}

himpunan G$_{\delta}$ {\bf 264Xd}











\indexheader{$\scriptstyle\pmb{H}$}




















































\indexmedskip
\indexheader{$\scriptstyle\pmb{L}$}



$\ell^1$ (dalam $\ell^1(X)$) {\bf 242Xa}, 243Xl, 246Xd, 247Xc,
247Xd


$\ell^1$ ($=\ell^1(\Bbb N)$) 246Xc


$\ell^2$ {\bf 244Xn}, 282K, 282Xg


$\ell^p$ (dalam $\ell^p(X)$) {\bf 244Xn}


$\ell^{\infty}$ (dalam $\ell^{\infty}(X)$) {\bf 243Xl}, {\bf 281B}, 281D


$\ell^{\infty}$ ($=\ell^{\infty}(\Bbb N)$) 243Xl












$\eusm L^0$ (dalam $\eusm L^0(\mu)$) 121Xb,
  \S241 ({\bf 241A}), \S245, 253C, 253Ya;
  (dalam $\eusm L^0_{\Sigma}$) {\bf 241Yc}
  {\it lihat juga} $L^0$ ({\bf 241C}),
$\eusm L^0_{\Bbb C}$ ({\bf 241J})



$\eusm L^0_{\Bbb C}$ (dalam $\eusm L^0_{\Bbb C}(\mu)$) {\bf 241J}, 253L


$L^0$ (dalam $L^0(\mu)$) \S241 ({\bf 241A}), 242B, 242J, {\it 243A},
{\it 243B}, {\it 243D}, {\it 243Xe}, 243Xj,
\S245, 253Xe-253Xg, 
271De, {\it 272H}; 


----- (dalam $L^0_{\Bbb C}(\mu)$) {\bf 241J}






----- {\it lihat juga} $\eusm L^0$ ({\bf 241A})


$\eusm L^1$ (dalam $\eusm L^1(\mu)$) {\it 122Xc}, 242A, 242Da, 242Pa,
242Xb; 
  (dalam $\eusm L^1_{\Sigma})$ {\bf 242Yg};
  (dalam $\eusm L^1_{\Bbb C}(\mu)$ {\bf 242P}, {\it 255Yn};
  (dalam $\eusm L^1_V(\mu)$) {\bf 253Yf};
  {\it lihat juga} $L^1$, $\|\,\|_1$
  

$L^1$ (dalam $L^1(\mu)$) \S242 ({\bf 242A}),
{\it 243De}, 243F, 243G, 243J, 243Xf-243Xh, 
{\it 243Xh},
245H, 245J, 245Xh, 245Xi, \S246,
\S247, \S253, 254R, 254Xq, 254Ya, 254Yd, 257Ya, {\it 282Bd} 


----- (dalam $L^1_V(\mu)$) {\bf 253Yf}, 253Yi




----- {\it lihat juga} $\eusm L^1$, $L^1_{\Bbb C}$, $\|\,\|_1$


$L^1_{\Bbb C}$ (dalam $L^1_{\Bbb C}(\mu)$) {\bf 242P}, 243K,
246K, {\it 246Yl}, {\it 247E}, 255Xi;
  {\it lihat juga} konvolusi fungsi


$\eusm L^2$ (dalam $\eusm L^2(\mu)$) 253Yj, \S286;  
  (dalam $\eusm L^2_{\Bbb C}(\mu)$) 284N, 284O, 284Wh, 284Wi, 284Xj,
284Xl-284Xn, 
284Yg;
  {\it lihat juga} $L^2$, $\eusm L^p$, $\|\,\|_2$


$L^2$ (dalam $L^2(\mu)$) 244N, 244Yl, {\it 247Xe}, 253Xe; 
  (dalam $L^2_{\Bbb C}(\mu)$) 244Pe, 282K, 282Xg, 284P
  {\it lihat juga} $\eusm L^2$, $L^p$, $\|\,\|_2$



$\eusm L^p$ (dalam $\eusm L^p(\mu)$) {\bf 244A}, 244Da, 244Eb, 244Pa,
244Xa, 244Ya, 244Yi,
246Xg, 252Yh, 253Xh, 255K, {\it 255Of}, 255Ye, 255Yf, 255Yk, 255Yl,
261Xa, 263Xa, 273M, 273Nb, 281Xd, 282Yc, 284Xk, 286A; 
  {\it lihat juga} $L^p$, $\eusm L^2$, $\|\,\|_p$


$L^p$ (dalam $L^p(\mu)$, $1<p<\infty$) \S244 ({\bf 244A}), 245G, 245Xk,
245Xl, 245Yg, 246Xh, 247Ya, 253Xe, 253Xi, 253Yk, 255Yh; 
  
  
  
  {\it lihat juga} $\eusm L^p$, $\|\,\|_p$


$\eusm L^{\infty}$ (dalam $\eusm L^{\infty}(\mu)$) {\bf 243A}, 243D,
{\it 243I}, 243Xa, 243Xl, 243Xn;
  
  {\it lihat juga} $L^{\infty}$


$\eusm L^{\infty}_{\Bbb C}$ {\bf 243K}


$\eusm L^{\infty}_{\Sigma}$ {\bf 243Xb}


$L^{\infty}$ (dalam $L^{\infty}(\mu)$) \S243 ({\bf 243A}),
253Yd


----- (dalam $L^{\infty}_{\Bbb C}(\mu)$) {\bf 243K}, 243Xm






-----  {\it lihat juga} $\eusm L^{\infty}$,
$\|\,\|_{\infty}$




$\eurm L$ (dalam $\eurm L(U;V)$, ruang dari operator linear) 253A,
253Xa








$\lim$ (dalam $\lim\Cal F$) {\bf 2A3S};
  (dalam $\lim_{x\to\Cal F}$) {\bf 2A3S}


$\liminf$ (dalam $\liminf_{n\to\infty}$) \S1A3 ({\bf 1A3Aa}), 2A3Sg;
  (dalam $\liminf_{t\downarrow 0}$) {\bf 1A3D}, 2A3Sg;
  (dalam $\liminf_{x\to\Cal F}$) {\bf 2A3S}
  

$\limsup$ (dalam $\limsup_{n\to\infty}$) \S1A3 ({\bf 1A3Aa}), 2A3Sg;
  (dalam $\limsup_{t\downarrow 0}$) {\bf 1A3D}, 2A3Sg;
  (dalam $\limsup_{x\to\Cal F}f(x)$) {\bf 2A3S}







$\ln^+$ {\bf 275Ye}


\indexmedskip
\indexheader{$\scriptstyle\pmb{M}$}









$\eusm M^{0,\infty}$ {\bf 252Yo}






$M^{1,\infty}$ (dalam $M^{1,\infty}(\mu)$) {\bf 244Xl}, 244Xm, 244Xo,
244Yd


















































$\med$ (dalam $\med(u,v,w)$) {\it lihat} fungsi median ({\bf
2A1Ac})




\indexmedskip
\indexheader{$\scriptstyle\pmb{N}$}

$\Bbb N$
{\it lihat} himpunan kuasa


$\Bbb N\times\Bbb N$ 111Fb























\indexmedskip
\indexheader{$\scriptstyle\pmb{P}$}

$\Cal P$ {\it lihat} himpunan kuasa



























p.p.\ (`presque partout') {\bf 112De}

$\Pr(X>a)$, $\Pr(\pmb{X}\in E)$ dan seterusnya\ {\bf 271Ad}


\indexmedskip
\indexheader{$\scriptstyle\pmb{Q}$}

$\Bbb Q$ (himpunan bilangan rasional) 111Eb, 1A1Ef





\indexmedskip
\indexheader{$\scriptstyle\pmb{R}$}

$\Bbb R$ (himpunan bilangan real) 111Fe, 1A1Ha,
  2A1Ha, 2A1Lb


$\Bbb R^I$ 245Xa, 256Yf;  
  
  {\it lihat juga} metrik Euklides, topologi Euklides




$\overline{\Bbb R}$ {\it lihat} garis real diperluas (\S135)



$\RoverC$ 2A4A


















\indexmedskip
\indexheader{$\scriptstyle\pmb{S}$}

$S$ (dalam $S(\frak A)$) 243I; 
  (dalam $S^f\cong S(\frak A^f)$) {\it 242M}, {\it 244Ha}




$\eusm S$ {\it lihat} fungsi uji menurun cepat ({\bf 284A})




$S^1$ (lingkaran satuan sebagai grup topologis) {\it lihat} grup lingkaran


$S^{r-1}$ (bola satuan dalam $\BbbR^r$) {\it lihat} sfera









$_{\text{sf}}$ (dalam $\mu_{\text{sf}}$)
{\it lihat} versi semifinit suatu ukuran ({\bf 213Xc});
  (dalam $\mu^*_{\text{sf}}$) {\bf 213Xg}, 213Xi


















\indexmedskip
\indexheader{$\scriptstyle\pmb{T}$}









$\Cal T$ (dalam $\Cal T_{\bar\mu,\bar\nu}$) 244Xm, 244Xo, 244Yd,
246Yc












$\frak T_m$ {\it lihat} konvergensi dalam ukuran ({\bf 245A})






$\frak T_s$ (dalam $\frak T_s(U,V)$)  
  {\it lihat} topologi lemah ({\bf 2A5Ia}),
  topologi lemah* ({\bf 2A5Ig})




\indexmedskip
\indexheader{$\scriptstyle\pmb{U}$}

$U$ (dalam $U(x,\delta)$) {\bf 1A2A}












\indexmedskip
\indexheader{$\scriptstyle\pmb{V}$}

$\Var$ (dalam $\Var(X)$) {\it lihat} varians ({\bf 271Ac});
  (dalam $\Var_Df$, $\Var f$) {\it lihat} variasi ({\bf 224A})




\indexmedskip
\indexheader{$\scriptstyle\pmb{W}$}



$w^*$-topologi {\it lihat} topologi lemah* ({\bf 2A5Ig})











\indexmedskip
\indexheader{$\scriptstyle\pmb{Z}$}

$\Bbb Z$ (himpunan bilangan bulat) 111Eb, 1A1Ee;
  (sebagai grup topologis) 255Xk
  







ZFC {\it lihat} teori himpunan Zermelo-Fraenkel

 



\indexmedskip\indexheader{$\scriptstyle\pmb{\beta}$} 

$\beta_r$ (volume dari bola satuan dalam $\BbbR^r$) 252Q, 252Xi,
{\it 265F}, {\it 265H}, {\it 265Xa}, {\it 265Xb}, {\it 265Xe}





\indexmedskip\indexheader{$\scriptstyle\pmb{\Gamma}$} 

$\Gamma$ (dalam $\Gamma(z)$) {\it lihat} fungsi gamma ({\bf 225Xh})



\indexmedskip\indexheader{$\scriptstyle\pmb{\Delta}$}





sistem-$\Delta$ {\it 2A1Pa}



















\indexmedskip\indexheader{$\scriptstyle\pmb{\lambda}$}



\indexmedskip\indexheader{$\scriptstyle\pmb{\mu}$}

$\mu_G$ (distribusi normal baku) {\bf 274Aa}




\indexmedskip\indexheader{$\scriptstyle\pmb{\nu}$}

$\nu_{\pmb{X}}$ {\it lihat} distribusi variabel acak
({\bf 271C})




\indexmedskip
\indexheader{$\scriptstyle\pmb{\pi}$}















Teorema $\pi$-$\lambda$ {\it lihat} Teorema Kelas Monoton (136B)

\indexmedskip





\indexmedskip
\indexheader{$\scriptstyle\pmb{\sigma}$}

aditif-$\sigma$ {\it lihat} aditif terhitung ({\bf 231C})






$\sigma$ himpunan {\bf 111A}, 111B,
111D-111G, 
111Xc-111Xf, 
111Yb, 136Xb, 136Xi,
  212Xh; 
  {\it lihat juga}
aljabar-$\sigma$ ({\bf 111G})


aljabar-$\sigma$ yang ditentukan oleh variabel acak {\bf 272C},
272D








lengkap-$\sigma$ {\it lihat} lengkap-$\sigma$ Dedekind
({\bf 241Fb})






medan-$\sigma$ {\it lihat} aljabar-$\sigma$ ({\bf 111A})







ruang ukur hingga-$\sigma$ {\bf 211D}, 211L, 211M, 211Xe,
{\it 212Ga}, {\it 213Ha}, 213Ma, 214Ia, 214Ka, 215B, 215C, 215Xe,
215Ya, 215Yb, {\it 216A}, {\it 232B}, {\it 232F},
234B, 234Ne, 234Xe, 234Xn, 235M, 235P,
235Xj, 241Yd, {\it 243Xi}, 245Eb, 245K, 245L, {\it 245Xe}, 251K,
251L, 251Wg, 251Wp, 252B-252E, 
252H, 252P, 252R, 252Xd, {\it 252Yb}, 252Yg, 252Yv







\indexheader{ideal-$\sigma$}


{$\sigma$ }
(himpunan) {\bf 112Db},
  211Xc, 212Xe, 212Xh


























\indexheader{subaljabar-$\sigma$}


subaljabar-$\sigma$ himpunan \S233 ({\bf 233A})


\indexheader{$\sigma$-subhomomorfisme}


\indexheader{$(\sigma,\infty)$-distributif}


\indexmedskip
\indexheader{$\scriptstyle\pmb{\Sigma}$}







$\sum_{i\in I}a_i$ {\bf 112Bd}, {\it 222Ba}, {\bf 226A}
  

\indexmedskip
\indexheader{$\scriptstyle\pmb{\tau}$}



\indexheader{aditif-$\scriptstyle\pmb{\tau}$}


ukuran aditif-$\tau$ {\it 256M}, 256Xb, 256Xc

























\indexmedskip
\indexiiheader{$\scriptstyle\pmb{\Phi}$}

$\Phi$ {\it lihat} fungsi distribusi normal ({\bf 274Aa})


\indexmedskip
\indexheader{$\scriptstyle\pmb{\chi}$}

$\chi$ (dalam $\chi A$, fungsi indikator dari suatu himpunan $A$)
{\bf 122Aa}





\indexmedskip
\indexheader{$\scriptstyle\pmb{\omega}$}

$\omega$ (ordinal tak hingga pertama) {\bf 2A1Fa}






$\omega_1$ (ordinal tak terhitung pertama) {\bf 2A1Fc}















$\omega_2$ (ordinal awal tak terhitung kedua)
{\bf 2A1Fc}














\indexmedskip

\indexiiheader{0}
$\{0,1\}^I$ (ukuran biasa pada) {\bf 254J}, 254Xd, 254Xe, 254Yd,
272N, 273Xb 




----- ----- (ketika $I=\Bbb N$) 254K, 254Xk, 254Xr, 256Xk, 256Yd,
261Yd 














$[0,1]^I$ (ukuran biasa pada) 254Ye


















\indexmedskip
\indexheader{simbol khusus}

$\setminus$ (dalam $E\setminus F$, `selisih himpunan') {\bf 111C}

$\symmdiff$ (dalam $E\symmdiff F$, `selisih simetris') {\bf 111C}







$\bigcup$ (dalam $\bigcup_{n\in\Bbb N}E_n$) {\bf 111C};
  (dalam $\bigcup\Cal A$) {\bf 1A1F}

$\bigcap$ (dalam $\bigcap_{n\in\Bbb N}E_n$) {\bf 111C};
  (dalam $\bigcap\Cal E$) {\bf 1A2F}

$\vee$, $\wedge$ (dalam suatu kisi) 121Xb, {\bf 2A1Ad}






$\restr$ (dalam $f\restr A$, pembatasan suatu fungsi pada suatu himpunan) {\bf
121Eh}

$\overline{\phantom{A}}$ {\it lihat} penutupan ({\bf 2A2A}, {\bf
2A3Db})


$\bar{\phantom{g}}$ (dalam $\bar h(u)$, dengan $h$ suatu fungsi Borel
dan $u\in L^0$) {\bf 241I}, 241Xd, 241Xi, 245Dd


$\eae$ {\bf 112Dg}, 112Xe, 241C

$\leae$ {\bf 112Dg}, 112Xe

$\geae$ {\bf 112Dg}, 112Xe







$\ll$ (dalam $\nu\ll\mu$) {\it lihat} kontinu mutlak ({\bf 232Aa})





















$*$ (dalam $f*g$, $u*v$, $\lambda*\nu$, $\nu*f$, $f*\nu$) {\it lihat}
konvolusi ({\bf 255E}, {\bf 255O}, {\bf 255Xh},
{\bf 255Xk}, {\bf 255Yn})


* (dalam topologi lemah*) {\it lihat} topologi lemah* ({\bf 2A5Ig}); 
 (dalam $U^*=\eurm B(U;\Bbb R)$, ruang topologis linear dual)
{\it lihat} dual ({\bf 2A4H});  
 
(dalam $\mu^*$) {\it lihat} ukuran luar yang diturunkan dari suatu ukuran ({\bf 132B})

$_*$ (dalam $\mu_*$) {\it lihat} ukuran dalam yang diturunkan dari suatu ukuran
(113Yh)

$'$
(dalam $T'$) {\it lihat} operator adjoin















$\int$ (dalam $\int f$, $\int fd\mu$, $\int f(x)\mu(dx)$) {\bf 122E}, {\bf
122K}, {\bf 122M}, 122Nb;
  {\it lihat juga} integral atas, integral bawah ({\bf 133I})

-----  (dalam $\int u$) {\bf 242Ab}, 242B, 242D


----- (dalam $\int_Af$) {\bf 131D}, {\bf 214D}, 235Xe;  
  {\it lihat juga} ukuran subruang

----- (dalam $\int_Au$) {\bf 242Ac}


$\overline{\int}$ {\it lihat} integral atas ({\bf 133I})

$\underline{\int}$ {\it lihat} integral bawah ({\bf 133I})





$\Rint$ {\it lihat} integral Riemann ({\bf 134K})



$|\,\,|$ (dalam suatu ruang Riesz) 241Ee, {\bf 242G}




$\|\,\|_1$ (pada $L^1(\mu)$) \S242 ({\bf 242D}), 246F, 253E, 275Xd,
282Ye;
  (pada $\eusm L^1(\mu)$) {\bf 242D}, 242Yg, 273Na, 273Xk


$\|\,\|_2$ {\bf 244Da}, 273Xl, 282Yf;
  {\it lihat juga} $L^2$, $\|\,\|_p$


$\|\,\|_p$ (untuk $1<p<\infty$) \S244 ({\bf 244Da}), 245Xj, 246Xb,
246Xh, 246Xi, 252Yh, 252Yo, 253Xe, 253Xh,
273M, 273Nb, 275Xe, 275Xf, 275Xh, 276Ya;
  {\it lihat juga} $\eusm L^p$, $L^p$




$\|\,\|_{\infty}$ {\bf 243D}, {\bf 243Xb}, {\bf 243Xo}, 244Xg,
273Xm, 281B;

  {\it lihat juga} supremum esensial ({\bf 243D}), $L^{\infty}$, $\eusm L^{\infty}$,
$\ell^{\infty}$








$\otimes$ (dalam $f\otimes g$) {\bf 253B}, 253C, 253I, 253J, 253L,
253Ya, 253Yb;
  (dalam $u\otimes v$) {\bf 253E}, 253F, 253G, 253L,
253Xc-253Xg, 
253Xi, 253Yd




$\tensorhat$ (dalam $\Sigma\tensorhat\Tau$) {\bf 251D}, 251K, 251M,
251Xa, 251Xl, 251Ya, 252P, 252Xe, 252Xh, 253C


$\Tensorhat$ (dalam $\Tensorhat_{i\in I}\Sigma_i$) {\bf 251Wb}, 251Wf,
{\bf 254E}, 254F, 254Mc, 254Xc, 254Xi, 254Xj



$\prod$ (dalam $\prod_{i\in I}\alpha_i$) {\bf 254F};
  (dalam $\prod_{i\in I}X_i$) {\bf 254Aa}


$\#$ (dalam $\#(X)$, kardinal dari $X$) {\bf 2A1Kb}






$\varsphat$, $\varspcheck$ (dalam $\varhatf$, $\varcheckf$) {\it lihat}
transformasi Fourier, transformasi Fourier invers ({\bf 283A})
















$\Reverse{\phantom{x}}$ (dalam $\Reverse{f}$) {\bf 284If}








$^+$
  (dalam $\kappa^+$, kardinal penerus) {\bf 2A1Fc}; 
  
  (dalam $f^+$, dengan $f$ suatu fungsi) {\bf 121Xa}, {\bf 241Ef}; 
  (dalam $u^+$, dengan $u$ termasuk dalam suatu ruang Riesz) {\bf 241Ef}; 
  
  (dalam $F(x^+)$, dengan $F$ suatu fungsi bernilai real) {\bf 226Bb}


$^-$ (dalam $f^-$, dengan $f$ suatu fungsi) {\bf 121Xa}, {\bf 241Ef};
  (dalam $u^-$, dalam suatu ruang Riesz) {\bf 241Ef};
  (dalam $F(x^-)$, dengan $F$ suatu fungsi bernilai real) {\bf 226Bb}
  






















$\infty$ {\it lihat} tak hingga





$[\hskip1em]$ (dalam $[a,b]$) {\it lihat} selang tertutup ({\bf 115G},
{\bf 1A1A}, {\bf 2A1Ab});
  (dalam $f[A]$, $f^{-1}[B]$, $R[A]$, $R^{-1}[B]$) {\bf 1A1B}

$[[\hskip1em]]$ (dalam $f[[\Cal F]]$) {\it lihat} filter citra
({\bf 2A1Ib})




$\coint{\hskip1em}$ (dalam $\coint{a,b}$) {\it lihat} selang setengah-terbuka
({\bf 115Ab}, {\bf 1A1A})

$\ocint{\hskip1em}$ (dalam $\ocint{a,b}$) {\it lihat} selang setengah-terbuka
({\bf 1A1A})

$\ooint{\hskip1em}$ (dalam $\ooint{a,b}$) {\it lihat} selang terbuka
({\bf 115G}, {\bf 1A1A})





$\fraction{\,}$ (dalam $\fraction{x}$, bagian pecahan) {\bf 281M}




$\LLcorner$ (dalam $\mu\LLcorner E$) {\bf 234M}, 235Xe








\frnewpage

